\documentclass[oneside,openany,phd]{TMU-Thesis}

% مشخصات پایان‌نامه را در فایلهای faTitle و enTitle وارد نمایید.
% در اینجا برای یکپارچگی، همینجا وارد می‌کنیم

% مشخصات فارسی
\university{تربیت مدرس}
\faculty{مهندسی برق و کامپیوتر}
\department{گروه هوش مصنوعی و رباتیکز}
\subject{مهندسی کامپیوتر}
\field{نرم‌افزار}
\credits{---} % در رساله دکتری معمولا به صورت واحدی ذکر نمی‌شود یا فیلد خالی است
\title{یادگیری تقویتی عمیق با شکل‌دهی تصادفی پاداش برای مدیریت عدم قطعیت در سیستم‌های خودوفقی}
\firstsupervisor{دکتر ...}
% \secondsupervisor{استاد راهنمای دوم}
\firstadvisor{دکتر ...}
% \secondadvisor{استاد مشاور دوم}
\name{علیرضا}
\surname{کاویانی‌فر}
\studentID{...}
\thesisdate{بهمن ۱۴۰۴}

% مشخصات انگلیسی
\latinuniversity{Tarbiat Modares University}
\latinfaculty{Faculty of Electrical and Computer Engineering}
\latinsubject{Computer Engineering}
\latinfield{Software}
\latintitle{Randomized Reward-Shaped Deep Reinforcement Learning for Managing Uncertainty in Self-Adaptive Systems}
\firstlatinsupervisor{Dr. ...}
\firstlatinadvisor{Dr. ...}
\latinname{Alireza}
\latinsurname{Kavianifar}
\latinthesisdate{February 2026}
\latinkeywords{Self-Adaptive Systems, Deep Reinforcement Learning, Reward Shaping, Uncertainty}

% فراخوانی بسته‌ها و دستورات استاندارد
% در این فایل، دستورها و تنظیمات مورد نیاز، آورده شده است.
%-------------------------------------------------------------------------------------------------------------------

% در ورژن جدید زی‌پرشین برای تایپ متن‌های ریاضی، این سه بسته، حتماً باید فراخوانی شود
\usepackage{amsthm,amssymb,amsmath}
% بسته‌ای برای تنطیم حاشیه‌های بالا، پایین، چپ و راست صفحه
\usepackage[top=30mm, bottom=25mm, left=25mm, right=35mm]{geometry}
% بسته‌‌ای برای ظاهر شدن شکل‌ها و تصاویر متن
\usepackage{graphicx}
% بسته‌ای برای رسم کادر
\usepackage{framed} 
% بسته‌‌ای برای چاپ شدن خودکار تعداد صفحات در صفحه «معرفی پایان‌نامه»
\usepackage{lastpage}
% بسته‌ و دستوراتی برای ایجاد لینک‌های رنگی با امکان جهش
% \usepackage[pagebackref=false,colorlinks,linkcolor=black,citecolor=black]{hyperref}
% چنانچه قصد پرینت گرفتن نوشته خود را دارید، خط بالا را غیرفعال و  از دستور زیر استفاده کنید چون در صورت استفاده از دستور زیر‌‌، 
% لینک‌ها به رنگ سیاه ظاهر خواهند شد که برای پرینت گرفتن، مناسب‌تر است
% \usepackage[pagebackref=false]{hyperref}
% بسته‌ لازم برای تنظیم سربرگ‌ها
\usepackage{fancyhdr}
\setlength{\headheight}{16pt}
%
\usepackage{setspace}
\usepackage{algorithm}
\usepackage{algorithmic}
\usepackage{subfigure}
\usepackage{titletoc}


% بسته‌ای برای ظاهر شدن «مراجع» و «نمایه» در فهرست مطالب
\usepackage[nottoc]{tocbibind}
% دستورات مربوط به ایجاد نمایه
\usepackage{makeidx}
\makeindex
% تعریف شیم‌های مورد نیاز برای نسخه‌های جدید بسته bidi در برخی محیط‌ها
% Removal of redundant and potentially broken definitions
\makeatletter
% شیم‌های سازگاری برای نسخه‌های جدید لاتک و بسته‌های bidi/xepersian
\providecommand{\IfClassLoadedT}[2]{\@ifclassloaded{#1}{#2}{}}
\providecommand{\IfClassLoadedF}[2]{\@ifclassloaded{#1}{}{#2}}
\providecommand{\IfClassLoadedTF}[3]{\@ifclassloaded{#1}{#2}{#3}}
\providecommand{\IfPackageLoadedT}[2]{\@ifpackageloaded{#1}{#2}{}}
\providecommand{\IfPackageLoadedF}[2]{\@ifpackageloaded{#1}{}{#2}}
\providecommand{\IfPackageLoadedTF}[3]{\@ifpackageloaded{#1}{#2}{#3}}
\providecommand{\@iffileloaded}[3]{\expandafter\ifx\csname ver@#1\endcsname\relax#3\else#2\fi}
\providecommand{\IfFileLoadedT}[2]{\@iffileloaded{#1}{#2}{}}
\providecommand{\IfFileLoadedF}[2]{\@iffileloaded{#1}{}{#2}}
\providecommand{\IfFileLoadedTF}[3]{\@iffileloaded{#1}{#2}{#3}}
\makeatother

%%%%%%%%%%%%%%%%%%%%%%%%%%
% فراخوانی بسته زی‌پرشین و تعریف قلم فارسی و انگلیسی
\usepackage[pagebackref=false,colorlinks,linkcolor=black,citecolor=black]{hyperref}
\usepackage{notoccite}

\ExplSyntaxOn
\cs_if_exist:NTF \g__hyp_linknestlevel_int { } { \int_new:N \g__hyp_linknestlevel_int }
\cs_if_exist:NTF \l__hyp_annot_GoTo_bool { } { \bool_new:N \l__hyp_annot_GoTo_bool }
\cs_if_exist:NTF \l__hyp_annot_URI_bool { } { \bool_new:N \l__hyp_annot_URI_bool }
\cs_if_exist:NTF \l__hyp_href_url_ismap_bool { } { \bool_new:N \l__hyp_href_url_ismap_bool }
\cs_if_exist:NTF \l__hyp_text_enc_uri_print_tl { } { \tl_new:N \l__hyp_text_enc_uri_print_tl }
\cs_if_exist:NTF \l__hyp_uri_tmpa_tl { } { \tl_new:N \l__hyp_uri_tmpa_tl }

% Defining missing l3pdf/pdfdict commands if not present
\cs_if_exist:NTF \pdfdict_put:nno { } { \cs_new:Npn \pdfdict_put:nno #1#2#3 { } }
\cs_if_exist:NTF \pdfdict_put:nnn { } { \cs_new:Npn \pdfdict_put:nnn #1#2#3 { } }
\cs_if_exist:NTF \pdfdict_use:n { } { \cs_new:Npn \pdfdict_use:n #1 { } }
\cs_if_exist:NTF \pdfannot_dict_put:nne { } { \cs_new:Npn \pdfannot_dict_put:nne #1#2#3 { } }
\cs_if_exist:NTF \pdfannot_link:nen { } { \cs_new:Npn \pdfannot_link:nen #1#2#3 { } }
\cs_if_exist:NTF \__hyp_check_link_nesting:TF { } { \cs_new:Npn \__hyp_check_link_nesting:TF #1#2 { #1 } }
\cs_if_exist:NTF \__hyp_text_pdfstring:eoN { } { \cs_new:Npn \__hyp_text_pdfstring:eoN #1#2#3 { } }
\ExplSyntaxOff

% تعریف دستورات گمشده برای جلوگیری از خطای renewcommand در بسته‌های داخلی
\makeatletter
\providecommand{\foottextfont}{}
\providecommand{\LTRfoottextfont}{}
\providecommand{\RTLfoottextfont}{}
\providecommand{\abstractname}{}
\providecommand{\latinabstract}{}
\providecommand{\AutoPersianDigits}{}
\makeatother

\usepackage{xepersian}
\settextfont[Path=font/,Extension=.ttf,Scale=1,
    BoldFont=XB NiloofarBd,
    ItalicFont=XB NiloofarIt,
    BoldItalicFont=XB NiloofarBdIt
]{XB Niloofar}
\setlatintextfont[Scale=0.9]{Times New Roman}

%%%%%%%%%%%%%%%%%%%%%%%%%%
% چنانچه می‌خواهید اعداد در فرمول‌ها، انگلیسی باشد، خط زیر را غیرفعال کنید
\setdigitfont[Path=font/,Extension=.ttf,Scale=1,
    BoldFont=XB NiloofarBd,
    ItalicFont=XB NiloofarIt,
    BoldItalicFont=XB NiloofarBdIt
]{XB Niloofar}
%%%%%%%%%%%%%%%%%%%%%%%%%%
% تعریف قلم‌های فارسی و انگلیسی اضافی برای استفاده در بعضی از قسمت‌های متن
\defpersianfont\titlefont[Path=font/,Extension=.ttf,Scale=1]{XB Titre}
% \defpersianfont\iranic[Scale=1.1]{XB Zar Oblique}%Italic}%
% \defpersianfont\nastaliq[Scale=1.2]{IranNastaliq}

%%%%%%%%%%%%%%%%%%%%%%%%%%
% دستوری برای حذف کلمه «چکیده»
\renewcommand{\abstractname}{}
% دستوری برای حذف کلمه «abstract»
%\renewcommand{\latinabstract}{}
% دستوری برای تغییر نام کلمه «اثبات» به «برهان»
\renewcommand\proofname{\textbf{برهان}}
% دستوری برای تغییر نام کلمه «کتاب‌نامه» به «مراجع»
\renewcommand{\bibname}{مراجع}
% دستوری برای تعریف واژه‌نامه انگلیسی به فارسی
\newcommand\persiangloss[2]{#1\dotfill\lr{#2}\\}
% دستوری برای تعریف واژه‌نامه فارسی به انگلیسی 
\newcommand\englishgloss[2]{#2\dotfill\lr{#1}\\}
% تعریف دستور جدید «\پ» برای خلاصه‌نویسی جهت نوشتن عبارت «پروژه/پایان‌نامه/رساله»
\newcommand{\پ}{پروژه/پایان‌نامه/رساله }

%\newcommand\BackSlash{\char`\\}

%%%%%%%%%%%%%%%%%%%%%%%%%%
\SepMark{-}

% تعریف و نحوه ظاهر شدن عنوان قضیه‌ها، تعریف‌ها، مثال‌ها و ...
\theoremstyle{definition}
\newtheorem{definition}{تعریف}[section]
\theoremstyle{plain}
\newtheorem{theorem}[definition]{قضیه}
\newtheorem{lemma}[definition]{لم}
\newtheorem{proposition}[definition]{گزاره}
\newtheorem{corollary}[definition]{نتیجه}
\newtheorem{remark}[definition]{ملاحظه}
\theoremstyle{definition}
\newtheorem{example}[definition]{مثال}

%\renewcommand{\theequation}{\thechapter-\arabic{equation}}
%\def\bibname{مراجع}
\numberwithin{algorithm}{chapter}
\def\listalgorithmname{فهرست الگوریتم‌ها}
\def\listfigurename{فهرست تصاویر}
\def\listtablename{فهرست جداول}

%%%%%%%%%%%%%%%%%%%%%%%%%%%%
% دستورهایی برای سفارشی کردن سربرگ صفحات
% \newcommand{\SetHeader}{
% \csname@twosidetrue\endcsname
% \pagestyle{fancy}
% \fancyhf{} 
% \fancyhead[OL,EL]{\thepage}
% \fancyhead[OR]{\small\rightmark}
% \fancyhead[ER]{\small\leftmark}
% \renewcommand{\chaptermark}[1]{%
% \markboth{\thechapter-\ #1}{}}
% }
%%%%%%%%%%%%5
%\def\MATtextbaseline{1.5}
%\renewcommand{\baselinestretch}{\MATtextbaseline}
\doublespacing
%%%%%%%%%%%%%%%%%%%%%%%%%%%%%
% دستوراتی برای اضافه کردن کلمه «فصل» در فهرست مطالب با استفاده از titletoc
\titlecontents{chapter}
  [5em]
  {\addvspace{1em}\bfseries}
  {\contentslabel[\chaptername~\thecontentslabel:]{5em}}
  {\hspace*{-5em}}
  {\titlerule*[0.5pc]{.}\contentspage}

\titlecontents{section}
  [5em]
  {}
  {\contentslabel{2em}}
  {\hspace*{-2em}}
  {\titlerule*[0.5pc]{.}\contentspage}

\titlecontents{subsection}
  [8em]
  {}
  {\contentslabel{3em}}
  {\hspace*{-3em}}
  {\titlerule*[0.5pc]{.}\contentspage}

\titlecontents{subsubsection}
  [12em]
  {}
  {\contentslabel{4em}}
  {\hspace*{-4em}}
  {\titlerule*[0.5pc]{.}\contentspage}

\makeatletter
{\def\thebibliography#1{\chapter*{\refname\@mkboth
   {\uppercase{\refname}}{\uppercase{\refname}}}\list
   {[\arabic{enumi}]}{\settowidth\labelwidth{[#1]}
   \rightmargin\labelwidth
   \advance\rightmargin\labelsep
   \advance\rightmargin\bibindent
   \itemindent -\bibindent
   \listparindent \itemindent
   \parsep \z@
   \usecounter{enumi}}
   \def\newblock{}
   \sloppy
   \sfcode`\.=1000\relax}}
\makeatother


% تنظیمات اختصاصی برای فهرست مطالب (TOC)
\setcounter{secnumdepth}{3}
\setcounter{tocdepth}{3}

% بازگرداندن نقاط رابط (Dot Leaders) که در Cleanup حذف شده بود
\renewcommand{\cftchapleader}{\cftdotfill{\cftdotsep}}
\renewcommand{\cftsecleader}{\cftdotfill{\cftdotsep}}
\renewcommand{\cftsubsecleader}{\cftdotfill{\cftdotsep}}

% تنظیمات تورفتگی و ظاهر
\setlength{\cftsubsubsecindent}{6em}
\setlength{\cftaftertoctitleskip}{10pt} % رفع فاصله زیاد بعد از عنوان فهرست
\raggedbottom

\begin{document}
\AutoPersianDigits
\setstretch{1.5}

\pagenumbering{harfi}
\firstPage
\besmPage
\davaranPage
% در این قسمت اسامی‌داوران در فایل اصلی باید دستی وارد شوند، اینجا از ساختار پیش‌فرض استفاده می‌کنیم

\esalatPage
\mojavezPage

\chapter*{تقدیم}
\addcontentsline{toc}{chapter}{تقدیم}
\begin{center}
    \textit{(متن تقدیم به خانواده، اساتید و ...)}
\end{center}
\newpage

\begin{acknowledgementpage}
(متن تشکر و قدردانی از اساتید راهنما، مشاور و دوستان)
\signature 
\end{acknowledgementpage}

\keywords{سیستم‌های خودوفقی، یادگیری تقویتی عمیق، شکل‌دهی پاداش، عدم قطعیت}
\fa-abstract{
(خلاصه‌ای از چالش، راهکار و نتایج کلیدی در ۳۰۰ تا ۵۰۰ کلمه)
}
\abstractPage

\tableofcontents
\newpage

\listoftables
\newpage

\listoffigures
\newpage

\chapter*{فهرست علائم و اختصارات}
\addcontentsline{toc}{chapter}{فهرست علائم و اختصارات}
(لیست نمادهای ریاضی مانند $\rho$، $\gamma$ و اختصارات مانند \lr{DRL}، \lr{SAS})
\newpage

% ============================================================================
% متن اصلی رساله
% ============================================================================
\pagenumbering{arabic}
\pagestyle{fancy}

% فصل ۱: مقدمه
\chapter{مقدمه}
\section{مقدمه}

در حوزه مهندسی سیستم‌های نرم‌افزاری معاصر، پیچیدگی به‌عنوان یکی از بارزترین ویژگی‌ها خودنمایی می‌کند. محیط عملیاتی پویا نقش محوری در شکل‌دهی به این پیچیدگی ایفا کرده و سیستم‌ها را مجبور می‌سازد تا با شرایط نامشخصی که معمولاً تنها پس از عملیاتی شدن بروز می‌کنند، دست و پنجه نرم کنند.

علاوه بر این، چالش‌ها و منابع عدم قطعیت در سیستم‌های مختلف به شکل‌های گوناگونی نمود پیدا می‌کنند. به‌عنوان مثال، در سیستم‌های مستقر در محیط‌های پویا نظیر برنامه‌های اینترنت اشیا (\lr{IoT}) در شهرهای هوشمند، تغییرات آب و هوایی تأثیر قابل توجهی می‌گذارد. همچنین محیط‌های ابری با تقاضای منابع متغیر، امنیت سایبری با تهدیدات نوظهور و مداوم، تجارت الکترونیک و شبکه‌های اجتماعی با رفتار غیرقابل پیش‌بینی کاربران، و سیستم‌های توزیع‌شده با تأخیر شبکه، خرابی سخت‌افزار و تغییرات زیرساختی روبه‌رو هستند.

سیستم‌های سنتی ایستا یا مبتنی بر قواعد، انعطاف‌پذیری لازم برای پاسخ به چنین تغییراتی در زمان اجرا را ندارند. در مقابل، سیستم‌های خودوفقی (\lr{SASs}) مجهز به حلقه‌های بازخوردی هستند که به آن‌ها اجازه می‌دهد تا محیط و وضعیت خود سیستم را پایش کنند، انحرافات از رفتار مورد انتظار را تحلیل نمایند و به‌طور خودمختار استراتژی‌ها و تغییرات وفقی را برنامه‌ریزی و اجرا کنند.

هنگامی‌که سیستم با یک فضای وفقی و گزینه‌های تغییر بسیار گسترده مواجه می‌شود، وظیفه انتخاب یک گزینه وفقی که با اهداف سیستم همسو باشد، از نظر محاسباتی بسیار پرهزینه و دشوار، و در برخی موارد تقریباً غیرممکن می‌گردد. تمرکز این رساله بر سیستم‌هایی است که دارای فضای وفقی گسسته و بزرگ هستند و نیاز فوری به راهکارهای کارآمد جهت انتخاب گزینه‌های بهینه در زمان اجرا دارند.

\section{بیان مسئله و سوالات تحقیق}

\subsection{مسئله تحقیق}

روش‌های موجود برای مدیریت عدم قطعیت در سیستم‌های خودوفقی با محدودیت‌های قابل توجهی روبه‌رو هستند. به‌عنوان مثال، روش‌های مبتنی بر یادگیری ماشین نیازمند داده‌های آموزشی برچسب‌دار هستند که نمایانگر محیط سیستم باشند؛ به دست آوردن چنین داده‌هایی در زمان طراحی به دلیل عدم قطعیت فراوان بسیار دشوار است. از طرفی، روش‌های مبتنی بر جستجو (مانند الگوریتم‌های صعود تپه و الگوریتم‌های ژنتیک) با مشکلات مقیاس‌پذیری مواجه هستند و عملکرد آن‌ها با تغییر شرایط محیطی دچار افت محسوس می‌شود.

یادگیری تقویتی عمیق (\lr{DRL}) اگرچه ابزاری قدرتمند است، اما خود با چالش‌های اساسی مانند گیر افتادن در بهینه‌های محلی مواجه می‌باشد. بهینه‌های محلی به راه‌حل‌هایی اشاره دارند که تنها در یک منطقه محدود از فضای مسئله بهینه هستند و به صورت سراسری بهترین راه‌حل ممکن محسوب نمی‌شوند. در زمینه \lr{DRL}، این بدان معناست که مدل ممکن است سیاستی را یاد بگیرد که تحت شرایط خاص به خوبی عملکرد مناسبی دارد، اما در صورت تغییر محیط قادر به تعمیم‌پذیری و حفظ عملکرد مطلوب نیست.

علاوه بر این، آموزش مجدد یا بازآموزی (\lr{Retraining}) مدل در زمان اجرا فرآیندی پرهزینه است. بازآموزی شامل مراحل متعددی مانند جمع‌آوری داده‌های جدید، ارزیابی مجدد ساختار و پارامترهای مدل، و تست مدل به‌روزشده برای اطمینان از عملکرد بهتر آن نسبت به نسخه قبلی است. در طول این فرآیند، سیستم ممکن است به صورت غیربهینه عمل کند که در سیستم‌های حیاتی، این دوره عملکرد نامطلوب می‌تواند پیامدهای جبران‌ناپذیری به همراه داشته باشد.

\textbf{بیان اصلی و خلاصه‌شده‌ی مسئله:} ماهیت پویا و غیرقابل پیش‌بینی محیط‌هایی که سیستم‌های نرم‌افزاری در آن‌ها عمل می‌کنند، چالش‌های قابل توجهی را برای مدل‌های سنتی یادگیری تقویتی عمیق به وجود می‌آورد. این مدل‌ها به‌طور مکرر با مشکل بهینه‌های محلی دست و پنجه نرم کرده و نمی‌توانند در زمان مواجهه با شرایط متغیر به خوبی عملکرد خود را تعمیم دهند که این امر منجر به عملکرد نامطلوب سیستم و نیاز به بازآموزی‌های مکرر و زمان‌بر خواهد شد.

\subsection{سوالات تحقیق}

با توجه به چالش‌ها و خلاءهای تحقیقاتی مطرح شده و بر اساس دستاوردهای مقاله پیشین \cite{RS_DRL_Article}، این رساله به دنبال پاسخگویی به سوالات اصلی زیر است:

\begin{enumerate}
    \item \textbf{سوال اول:} چگونه می‌توان یک روش مبتنی بر یادگیری تقویتی عمیق طراحی کرد که به‌طور موثری از بهینه‌های محلی در محیط‌های پویا، بدون نیاز به بازآموزی مکرر، فرار کند؟
    \item \textbf{سوال دوم:} چگونه می‌توان مکانیزم‌های شکل‌دهی پاداش (\lr{Reward Shaping}) را با بازپخش تجربه (\lr{Experience Replay}) ادغام کرد تا تعمیم‌پذیری سیستم در تغییرات شرایط محیطی بهبود یابد؟
    \item \textbf{سوال سوم:} چگونه روش پیشنهادی می‌تواند فضاهای وفقی گسسته و بزرگ (از ۲۱۶ تا ۴۰۹۶ گزینه) را در حین حفظ عملکرد نزدیک به بهینه، به شکل کارآمدی مدیریت کند؟
    \item \textbf{سوال چهارم:} تا چه حد روش پیشنهادی (\lr{RS-DRL}) از روش‌های موجود نظیر یادگیری ای‌گرایانه حریصانه (\lr{$\epsilon$-greedy DQN})، \lr{DLASER} و \lr{ML4EAS} از منظر عملکرد مجانبی، سرعت همگرایی، و بهینه‌سازی چندهدفه، برتری و عملکرد بهتری از خود نشان می‌دهد؟
\end{enumerate}

\section{اهداف رساله}

این رساله اهداف کلان و خرد زیر را در راستای توسعه تحقیقات پیشین \cite{RS_DRL_Article} دنبال می‌کند:

\begin{enumerate}
    \item \textbf{طراحی و توسعه:} طراحی یک معماری نوآورانه یادگیری تقویتی عمیق با شکل‌دهی پاداش تصادفی (\lr{RS-DRL}) که به شکل یکپارچه‌ای با حلقه خودمختار \lr{MAPE-K} در سیستم‌های خودوفقی ادغام شده باشد.
    \item \textbf{نوآوری الگوریتمی:} توسعه یک مکانیزم نوین شکل‌دهی پاداش که تجربیات ناموفق را در طول فرآیند بازپخش حافظه به صورت تصادفی بازطراحی می‌نماید تا اکتشاف (\lr{Exploration}) را ترویج کرده و از همگرایی مدل در بهینه‌های محلی جلوگیری کند.
    \item \textbf{مدل‌سازی صوری:} مدل‌سازی صوری مسئله کنترل و وفقی‌سازی تحت عدم قطعیت که شامل فضای حالت (نظیر سطح انرژی، از دست رفتن بسته‌های شبکه، زمان تأخیر)، فضای عمل (گزینه‌های وفقی گسسته) و توابع پاداش جهت بهینه‌سازی‌های تک‌هدفه و چندهدفه است.
    \item \textbf{ارزیابی تجربی:} ارزیابی راهکار توسعه‌یافته \lr{RS-DRL} در تقابل با روش‌های پیشرفته‌ی فعلی (مانند \lr{Soft HER}، \lr{HER}، \lr{$\epsilon$-greedy DQN}، \lr{DLASER} و \lr{ML4EAS}) بر روی مطالعات موردی سیستم \lr{DeltaIoT} با اندازه‌های مختلف از فضاهای تطبیق (از جمله ۲۱۶، ۱۲۹۶ و ۴۰۹۶ گزینه ممکن).
    \item \textbf{تحلیل‌های آماری:} ارائه ارزیابی‌های آماری دقیق (تست مستقل \lr{t-test}، مقادیر \lr{p-value} و فاصله‌های اطمینان) برای تصدیق برتری روش پیشنهادی.
\end{enumerate}

\section{نوآوری‌های رساله}

\subsection{نوآوری اول: ارائه الگوریتم یادگیری تقویتی عمیق با شکل‌دهی تصادفی پاداش (\lr{RS-DRL})}
برخلاف بازپخش تجربه سنتی در یادگیری تقویتی، الگوریتم پیشنهاد شده در این رساله (\lr{RS-DRL}) برخی از تجربیات ناموفق گذشته را بازطراحی کرده و با اعطای پاداش مثبت به آن‌ها (همانند پاداش‌های موفقیت)، فرآیند یادگیری انسانی از اشتباهات گذشته خود را شبیه‌سازی می‌کند. مکانیزم انتخاب تصادفی تجربیات ناموفق که به وسیله فاکتور شکل‌دهی $\rho$ کنترل می‌شود، تجربیات مردود سیستم را با تخصیص پاداش خوش‌بینانه (به‌عنوان مثال تخصیص $r_i \leftarrow 1$) شبیه‌سازی موفقیت می‌نماید. این ترفند برای محیط‌های متنوع نیازی به تعاریف تخصصی و دانش دامنه خاص ندارد و موجب می‌شود تا: ۱) تعمیم‌پذیری را در طول محیط‌های متنوع تشویق نماید و ۲) فرآیند یادگیری را از طریق تبدیل اشتباهات گذشته به درس‌های با ارزش در تجربه‌های بعدی، به شدت تسریع کند.
در مقایسه با روش‌هایی نظیر بازپخش تجربه با نگاه به گذشته (\lr{HER}) و \lr{Soft HER} که ملزم به داشتن دانش دامنه‌ای و داشتن هدف صریح برای برچسب‌زنی مجدد اهداف (\lr{Goal Relabeling}) هستند، رویکرد ما در \lr{RS-DRL} از استراتژی پاداش تصادفی و خوش‌بینانه و کاملاً مستقل از دانش حوزه‌ای بهره می‌برد.

\subsection{نوآوری دوم: یکپارچه‌سازی \lr{RS-DRL} با معماری خودوفقی \lr{MAPE-K}}
نوآوری استراتژیک دیگر، ارائه چارچوب معماری تلفیقی شامل سیستم مدیریتی و سیستم مدیریت‌شونده است. معماری پیشنهادشده که با حلقه کنترل و بازخورد \lr{MAPE-K} ادغام شده است، شامل چرخه ارزیابی و انتخاب ۹ مرحله‌ای در تعامل بین مؤلفه‌های مختلف آن در زمان اجرا خواهد بود.
این اجزای کلیدی شامل ماژول استخراج‌کننده ویژگی از داده‌های سنسورها، انتخابگر فضای گزینه (\lr{$\epsilon$-greedy DQN})، تایید‌کننده گزینه وفقی از طریق اعتبارسنجی مدل با موتور \lr{UPPAAL-SMC}، محاسبه‌گر پاداش، حافظه بازپخش با در نظر گرفتن مکانیزم بدیع شکل‌دهی پاداش و مؤلفه آموزش مدل می‌باشد.

\subsection{نوآوری سوم: مدیریت موثر عدم قطعیت در فضاهای وفقی بزرگ و گسسته}
به‌عنوان یک نوآوری اساسی در سطح ارزیابی و دستاوردهای سیستم، این راهکار نشان داده که پتانسیل چشم‌گیری در مدیریت و راهبری سیستم‌های با مقیاس‌پذیری فضاهای حالت و فضای عمل بالا، مانند شبکه‌های گسترده‌ی اینترنت اشیاء داراست.
کاربست رویکرد \lr{RS-DRL} در پلتفرم \lr{DeltaIoT} از طریق کاهش چشم‌گیر هدررفت داده‌ها تا محدوده ۲۲.۲۳٪ نسبت به رقبای سرسختی مانند \lr{DLASER+} و بهبود ۳۹.۰۷ درصدی زمان تأخیر انتقال در شبکه، قابل توجه است. استفاده از فضای گسسته وفقی مقیاس‌پذیر در محیط‌های آزمایشگاهی شامل نسخه‌ی \lr{DeltaIoTv1} با ۲۱۶ گزینه، نسخه‌ی \lr{DeltaIoTv1.1} شامل ۱۲۹۶ گزینه و نسخه‌ی \lr{DeltaIoTv2} شامل ۴۰۹۶ حالت ممکن، کارایی و میزان پاداش مجانبی بهبود یافته و پایداری را نسبت به روش‌های پیشین نشان می‌دهد.

\section{مروری بر فصول رساله}

ساختار رساله به گونه‌ای طراحی شده است که به‌طور نظام‌مند روند توسعه، اجرا و ارزیابی رویکرد پیشنهادی را در اختیار خواننده قرار دهد. در ادامه، مرور خلاصه‌ای از هر یک از فصول آورده شده است:

\textbf{فصل ۲: مبانی و پیشینه تحقیق}\\
مفاهیم بنیادی مربوط به سیستم‌های خودوفقی، چالش‌های عدم قطعیت و مبانی مدل‌سازی آن در زمان اجرا، اصول سیستم‌های یادگیری تقویتی (\lr{MDP}، یادگیری \lr{Q}) و \lr{DRL} (به ویژه الگوریتم‌های \lr{DQN}) به تفصیل شرح داده می‌شود. سپس کارهای مرتبط پیشین در مدیریت فضاهای بزرگ تطبیق، از روش‌های مبتنی بر یادگیری ماشین نظیر \lr{ML4EAS} گرفته تا الگوریتم‌های اکتشافی جستجو و سایر روش‌های نوین مورد بررسی و ارزیابی قرار می‌گیرند.

\textbf{فصل ۳: مدل سیستم و فرمول‌بندی مسئله}\\
در این فصل، سیستم شبکه‌ی سنسورهای اینترنت اشیاء (\lr{DeltaIoT}) به‌عنوان مطالعه‌ی موردی انگیزشی معرفی می‌شود. فضای حالت صوری شبکه شامل مفاهیمی‌پیرامون مصرف انرژی، از دست‌رفتگی بسته‌های داده و زمان تأخیر ($S = \{\text{\lr{Energy}}, \text{\lr{Packet Loss}}, \text{\lr{Latency}}\}$) فرمول‌بندی شده و در ادامه، گزینه‌های فضای اعمال گسسته نیز مشخص و تدوین می‌گردد. در نهایت، توابع پاداش مربوط به سناریوهای بهینه‌سازی چندهدفه و تک‌هدفه به صورت ریاضی ارائه می‌شوند.

\textbf{فصل ۴: روش پیشنهادی \lr{RS-DRL}: معماری و الگوریتم}\\
مرکز ثقل این رساله در این فصل واقع شده است که در آن، معماری پیشنهادی یکپارچه با کنترل‌کننده‌ی \lr{MAPE-K} تبیین می‌گردد. الگوریتم پایه (\lr{RS-DRL})، به انضمام جزئیات دقیق مکانیزم شکل‌دهی پاداش‌های تصادفی با جزئیات الگوریتمی‌تشریح گردیده و پیاده‌سازی‌های مرتبط با اعتبارسنجی استراتژی‌های اجرایی پیش از پیاده‌سازی نهاییِ تغییرات توسط نرم‌افزار \lr{UPPAAL-SMC}، برای تضمین پایداری سیستم، بحث و بررسی می‌شوند.

\textbf{فصل ۵: ارزیابی تجربی}\\
نحوه تخصیص محیط شبیه‌سازی، سناریوهای متعدد پیکربندی شبکه، و پروتکل‌های مقایسه‌ی جامع مدل \lr{RS-DRL} با استراتژی‌های به‌روزِ رقیب نظیر \lr{$\epsilon$-Greedy DQN}، \lr{HER}، \lr{Soft HER}، \lr{DLASER}، و \lr{ML4EAS} در این فصل ارائه می‌گردد. ارائه‌ی شفافِ معیارهای کلیدی دستاوردها شامل بررسی میزان ظرفیت عملکرد مجانبی سیستم، ارزیابی زمان مورد نیاز تقربی برای رسیدن به آستانه قابل اتکا، و در نهایت کل مقادیر پاداش دریافتی در هر دوره‌ی آزمایش همراه با تأییدیه‌های تحلیل آماری و علمی‌مستند می‌باشند.

\textbf{فصل ۶: بحث، کارهای آینده و نتیجه‌گیری}\\
استلزامات نظریِ مرتبط با علت تفوقِ راهکار پیشنهادی بازبینی گردیده و کارایی مکانیزم از زوایای پراتیک صنعتی بحث می‌شود. در ادامه‌ی این فصل، به بررسی مصالحه‌ها، محدودیت‌ها و چالش‌های باقی‌مانده‌ی این معماری پرداخته شده و در نهایت، مسیرها و چارچوب‌هایی الهام‌بخش جهت بسط رویکرد فعلی در سیستم‌های پیوسته‌ی سایبرفیزیکی آینده پیشنهاد می‌گردد. فصل با نتیجه‌گیری کلی پیرامون مزیت‌های کاربردی راهکار نوین برای سیستم‌های عملیاتی بسته خواهد شد.

\newpage

% فصل ۲: مبانی و پیشینه تحقیق
\chapter{مبانی و پیشینه تحقیق}
\section{مقدمه}

هدف این فصل، ایجاد و بسط یک چارچوب و زیربنای نظری مستحکم است تا خواننده را از مفاهیم سطح بالا به سمت درک دقیق شکاف تحقیقاتی که این رساله آن را پر می‌کند، هدایت نماید. سیستم‌های نرم‌افزاری مدرن همواره با عدم قطعیت‌های مداوم مواجه هستند و باید در زمان اجرا به‌طور پیوسته ساختار و رفتار خود را با شرایط در حال تغییر وفق دهند. این فصل در ابتدا اصول اولیه سیستم‌های خودوفقی و مؤلفه‌های آن را تشریح می‌کند، سپس به نقش یادگیری تقویتی و محدودیت‌های روش‌های سنتی در مقابله با فضاهای بزرگ تصمیم‌گیری می‌پردازد و در نهایت با مروری ساختاریافته بر ادبیات تحقیق، علت نیاز به راهکار نوین \lr{RS-DRL} را اثبات می‌نماید.

\section{مبانی سیستم‌های خودوفقی}
\subsection{مفاهیم اصلی و الگوهای معماری}

سیستم‌های خودوفقی (\lr{Self-Adaptive Systems}) یک رویکرد برجسته و مؤثر برای مقابله با عدم قطعیت به شمار می‌روند. ایده کلیدی در این سیستم‌ها، تجهیز آن‌ها به یک یا چند حلقه بازخوردی کنترلی است که به صورت مداوم مؤلفه‌های داخلی سیستم و محیط خارجی آن را پایش کرده و در صورت نیاز سیستم را با تغییرات منطبق می‌نمایند \cite{RS_DRL_Article}. 

مهمترین و شناخته‌شده‌ترین الگوی معماری برای پیاده‌سازی این سیستم‌ها، حلقه خودمختار \lr{MAPE-K} است. همان‌طور که در فصل‌های بعدی با جزئیات بیشتر بررسی خواهد شد، این حلقه در کنار سیستم هدف و منطق کسب‌وکار، جریان کنترلی زیر را اعمال می‌کند:
\begin{itemize}
    \item \textbf{پایش (\lr{Monitor}):} سنسورها وضعیت محیط فیزیکی، زیرساخت و خود سیستم را به صورت زمان‌حقیقی استخراج می‌کنند.
    \item \textbf{تحلیل (\lr{Analyze}):} داده‌های پایش‌شده را پردازش کرده و مشخص می‌کند که آیا انحرافی از اهداف کیفی و کمی‌سیستم رخ داده است و نیاز به اعمال تغییر وجود دارد یا خیر.
    \item \textbf{برنامه‌ریزی (\lr{Plan}):} در صورت نیاز تطبیق، راهکارها و استراتژی‌های موجود (گزینه‌های وفقی) مقایسه شده و گزینه بهینه برای شرایط فعلی برگزیده می‌شود.
    \item \textbf{اجرا (\lr{Execute}):} استراتژی برنامه‌ریزی‌شده در قالب یک سری از دستورات اجرایی از طریق عملگرها (Actuators) بر روی سیستم اعمال می‌گردد.
    \item \textbf{دانش (\lr{Knowledge}):} به‌عنوان حافظه مرکزی عمل کرده و اطلاعات، مدل‌ها و داده‌های بین چهار مرحله قبلی را هماهنگ و نگهداری می‌کند. راهکار پیشنهادی ما (\lr{RS-DRL}) مستقیماً به‌عنوان هسته اصلی بخش برنامه‌ریزی در این الگو فعالیت می‌کند.
\end{itemize}

\subsection{چالش‌های عدم قطعیت در زمان اجرا}

عدم قطعیت در سیستم‌های خودوفقی، به هرگونه انحراف از دانش قطعی اطلاق می‌شود که می‌تواند باعث کاهش اطمینان نسبت به تصمیمات زمان اجرا گردد \cite{RS_DRL_Article}. 

به‌عنوان نمونه‌های عینی، می‌توان به موارد گوناگونی در دامنه‌های مختلف اشاره کرد: در برنامه‌های مبتنی بر اینترنت اشیاء (\lr{IoT}) که در محیط‌های شهری پراکنده هستند، تغییرات مداوم و غیرقابل پیش‌بینی آب و هوا منجر به افت کیفیت سیگنال و اختلال در ارسال اطلاعات سنسورها می‌شود. در محیط‌های ابری، نوسان شدید در تقاضای منابع از سوی کاربران، و در حوزه‌های امنیتی، ظهور بی‌وقفه و تصادفی حملات سایبری از مصادیق برجسته عدم قطعیت هستند.

اهمیت و چالش‌برانگیز بودن عدم قطعیت از آنجا نشأت می‌گیرد که مفروضات، پیکربندی‌ها و تصمیمات در نظر گرفته شده در فاز طراحی معماری نرم‌افزار را بی‌اعتبار می‌سازد. در نتیجه، سیستم باید قادر باشد با یادگیری آنلاین و پویا در زمان اجرا، اهداف اصلی خود را محقق نماید؛ وگرنه با نقض کیفیت خدمات صدمات سنگین احتمالی را تجربه خواهد نمود.

\section{مدل‌سازی عدم قطعیت}
\subsection{منابع عدم قطعیت}

منابع عدم قطعیت به‌طور کلی در سه دسته وسیع طبقه‌بندی می‌شوند: نخست، نامشخص بودن موجودیت‌های محیطی نظیر پارازیت در حسگرها، نویزها و تغییرات جوی که از کنترل ساختار درونی سیستم خارج هستند. دوم، ماهیت غیرقابل اتکای منابع خود سیستم مثل نقص در قطعات زیربنایی، کاهش سریع سطح باتری یا مشکلات نرم‌افزاری تعاملات میان سرویس‌ها؛ و در نهایت عدم قطعیت در اهداف متغیر و ترجیحات نامشخص ذینفعان در محیط‌های تجاری.

\subsection{رویکردهای مدل‌سازی عدم قطعیت}

جهت درک و محاسبه رفتار سیستم در هنگام بروز مقادیر ناشناخته، استفاده از مدل‌سازی‌های زمان اجرا (\lr{Runtime Models}) ضروری به نظر می‌رسد. در شرایطی که فضای جستجو بسیار بزرگ و ریسک تصمیم‌گیری به دلیل شرایط مبهم بالاست، رویکرد پیشنهادی از مکانیزم‌های مدل‌سازی ریاضی صوری بهره‌گیری و پشتیبانی می‌کند.

در این راستا ابزارهایی مانند موتور \lr{UPPAAL-SMC} و چارچوب \lr{ActivFORMS} به کار گرفته می‌شوند. این زیرساخت‌ها با استفاده از اتوماتای زمان‌دار (\lr{Timed Automata}) مدلی رفتاری از سیستم هدف و محیط آن را توسعه داده و از طریق روش‌های کنترل مدل آماری (\lr{Statistical Model Checking}) می‌توانند به سرعت گزینه‌های وفقی پیشنهاد شده توسط سیستم هوشمند را پیش از اجرا در سیستم فیزیکی ارزیابی، تایید یا رد کنند. این قابلیت نقشی کلیدی در امنیت راهکار پیشنهادی (در بخش تاییدگر گزینه وفقی) این رساله ایفا می‌کند.

\section{یادگیری تقویتی برای وفقی‌سازی}
\subsection{اصول یادگیری تقویتی (\lr{MDP}، یادگیری \lr{Q})}

یادگیری تقویتی از طریق تعامل مستمر یک عامل (\lr{Agent}) با محیط اطرافش بر پایه آزمون و خطا بنا شده است و این محیط به لحاظ ریاضی از طریق فرآیند تصمیم‌گیری مارکوف (\lr{MDP}) فرموله می‌شود. فرآیند \lr{MDP} خود شامل فضای حالت (\lr{State})، فضای عمل (\lr{Action}) و توابع پاداش (\lr{Reward}) است.

در فرآیند تصمیم‌گیری مارکوف، برای هر حالت خاص، عامل عمل مشخصی را اتخاذ می‌کند و مطابق با معادله مشهور بلمن (\lr{Bellman Equation}) محیط ارزش عمل اتخاذ شده را در قالب یک پاداش ارزیابی می‌نماید و به وضعیت بعدی منتقل می‌شود. این مفهوم در سیستم‌های خودوفقی مانند محیط \lr{DeltaIoT} به این صورت اعمال می‌گردد که وضعیت سیستم با معیارهایی نظیر انرژی ذخیره شده سیستم و تاخیر اندازه‌گیری شده فضای حالت را تشکیل می‌دهند و اقدام عامل، همان گزینه وفقی مورد انتخاب خواهد بود و محاسبه پاداش نیز طبق معادله بهینه‌سازی سیستم حاصل می‌گردد. نقش الگوریتم مشهور یادگیری Q ایجاد جدولی جامع برای تخصیص و پیدا کردن ارزش‌هاست تا مدل به سیاست بهینه همگرا شود.

\subsection{یادگیری تقویتی عمیق (\lr{DQN}) برای فضاهای بزرگ}

اما زمانی که با سیستم‌های واقعی که دارای فضای عمل گسسته ولی بسیار پهناور (برای مثال ۴۰۹۶ گزینه وفقی در شبکه‌های اینترنت اشیاء) مواجه هستیم، استفاده از متد کلاسیک یادگیری \lr{Q} با چالش نفرین ابعاد (\lr{Curse of Dimensionality}) روبه‌رو می‌شود؛ زیرا ذخیره و بروزرسانی جدولی با میلیون‌ها حالت و عمل از لحاظ محاسباتی ناکارآمد و در بسیاری موارد غیرممکن است.

شبکه‌های \lr{Q} عمیق (\lr{DQN}) راهکاری هوشمندانه برای این محدودیت ارائه می‌دهند؛ آن‌ها به جای جداول، تابع ارزش کیفیت \lr{Q} را با استفاده از شبکه‌های عصبی عمیق تقریب می‌زنند. برای غلبه بر بی‌ثباتی آموزش شبکه‌های عصبی در \lr{RL}، مدل \lr{DQN} از دو نوآوری برجسته بهره می‌برد: نخست استفاده از مکانیسم حافظه بازپخش تجربه (\lr{Experience Replay}) که نمونه‌های آموزشی تصادفی‌سازی شده و باعث شکستن همبستگی زمانی متوالی می‌گردد؛ و دوم تعبیه هدف ثابت به واسطه شبکه هدف ثابت (\lr{Target Network}) که بهبود تثبیت توابع هزینه را در هنگام محاسبه پیش‌بینی‌ها به دنبال دارد.

\subsection{مفهوم یادگیری مداوم در حلقه‌های وفقی}

اصلی‌ترین کاربرد این رساله، ادغام مفهوم آموزش پیوسته و مداوم توسط عامل \lr{DQN} در چارچوب معماری خودوفقی و به‌طور خاص درون حلقه \lr{MAPE-K} (همانطور که در فصل ۴ و شکل ۴ از مقاله \cite{RS_DRL_Article} به تفصیل نشان داده خواهد شد) است. پس از جاگذاری در مؤلفه برنامه‌ریزی (\lr{Planner})، عامل یادگیرنده وضعیتی را به‌عنوان ورودی فضای حالت از ماژول پایش دریافت نموده، استراتژی اقدام را برنامه‌ریزی کرده و بر اساس میزان موفقیت انتخابش از سوی مؤلفه تحلیل پاداش می‌گیرد. این ساختار تعاملی، سیستم خودوفقی را از پیکربندی ایستا به ساختار یادگیری پویا تبدیل می‌سازد.

\section{کارهای مرتبط: مدیریت فضاهای وفقی بزرگ}

دستیابی به کارایی مناسب در مدیریت محیط‌های مبهم بسیار کلیدی است. برای تبیین عمیق این مباحث، در ادامه به مقایسه انتقادی دسته‌بندی‌های موجود می‌پردازیم.

\subsection{روش‌های مبتنی بر یادگیری ماشین (\lr{ML-based Methods})}

اولین دسته‌بندی کارهای مرتبط در به‌کارگیری یادگیری ماشین نظارت‌شده و پیش‌بینی‌کننده با هدف هرس و کاهش فضای جستجو متمرکز است. رویکرد محققانی چون کین (\lr{Quin}) و همکارانش \cite{quin2019efficient,quin2022reducing} و مدل یادگیری برای ارزیابی کیفی ماشین (\lr{ML4EAS}) توسعه یافته توسط کاچی و بواناکا \cite{kachi2023hybrid} نشان می‌دهند که الگوریتم‌های طبقه‌بندی و رگرسیون قادرند تعداد زیادی از گزینه‌های غیربهینه را تا پیش از ارسال به مؤلفه تحلیل‌گر، غربال کنند.

\textbf{نقد رویکرد:} اگرچه این روش‌ها در کاهش اندازه محیط عمل خود کارا هستند، اما وجود وابستگی بنیادین به مجموعه‌داده‌های با برچسب آموزش‌دیده، نقص متدولوژی است. تهیه مجموعه‌داده در زمان توسعه که همسان با متغیرهای غیرمترقبه واقعیت باشد غیرعملی است. متعاقباً با بروز رانش مفهومی‌(\lr{Concept Drift}) ناشی از تغییر مداوم شرایط محیطی، این مدل‌های پیش‌آموزش دیده، عملکرد تعمیم‌یافته و کارآمد خود را در برخورد با رویه‌های ناشناخته از دست می‌دهند.

\subsection{روش‌های مبتنی بر جستجو (\lr{Search-based Methods})}

متدولوژی کلاسیک دیگر برای پیمایش سریع فضاهای وسیع جستجو استفاده از راهبردهای اکتشافی است. برای مثال استفاده کوکر (\lr{Coker}) از الگوریتم صعود تپه (\lr{Hill Climbing}) \cite{coker2015sass} و برنامه‌ریزی‌های مبتنی بر الگوریتم ژنتیک (\lr{Genetic Programming}) که به قلمرو مقالات کینر (\lr{Kinneer}) و همکاران اختصاص دارد \cite{kinneer2018managing,kinneer2020building,kinneer2021information}. همچنین شبیه‌سازی جستجو با استفاده از رویکرد تکامل چند هدفه نظیر \lr{MOEA} توسط چن (\lr{Chen}) \cite{chen2018femosaa}، از کاربردهای متداول است.

\textbf{نقد رویکرد:} هرچند ویژگی انطباق خودکار بدون آموزش در این رویکرد به چشم می‌خورد، اما چالش مقیاس‌پذیری، مانع استقرار کامل آن‌ها در پروژه‌های بلادرنگ است. اجرای مجدد جمعیت تکاملی همچون ژنتیک برای یافتن سناریوی مطلوب در پی هر گونه تطابق زمانی، مستلزم محاسبات بی‌نهایت پرهزینه و ایجاد تاخیر زیاد است؛ موضوعی که در شبکه‌های حساس تاخیرپذیر غیرقابل تحمل است.

\subsection{روش‌های مبتنی بر یادگیری تقویتی (\lr{RL-based Methods})}

پیشرفت‌های حاصل شده در اتوماسیون برخط سیستم‌ها با کارهای کیم و پارک \cite{kim2009reinforcement}، متزگر \cite{metzger2022realizing} و مدل \lr{RLMC} توسط کامارا \cite{camara2020quantitative} به دلیل تعامل بی‌وقفه پویایی‌های آنلاین شناخته شده است که عمدتا براساس کلاسیک الگوریتم \lr{Q-Learning} توسعه یافتند.

\textbf{نقد رویکرد:} الگوریتم کلاسیک در ابعاد بزرگ، همانند شبکه‌ای با میلیاردها ماتریس برای مدل‌سازی از پای در می‌آید و ناتوان از رسیدگی است. حتی تلاش‌هایی که با مدل پایه ارزیابی ما یعنی معماری حریصانه کلاسیک \lr{DQN} با استراتژی $\epsilon$\lr{-greedy} صورت پذیرفته است، کماکان با مشکل شدید محدودیت روبه‌رو است؛ زیرا مدل پایه عامل هوشمندی را پرورش می‌دهد که اغلب در نقاط و قله‌های بهینه‌های محلی (\lr{Local Optima}) متوقف شده، مانع قابلیت یافتن قله‌های اصلی (\lr{Global Optima}) در دینامیک شرایط محیط می‌گردد.

\subsection{شکل‌دهی پاداش و بازپخش تجربه با نگاه به گذشته (\lr{Reward Shaping \& HER})}

بازپخش تجربه‌های ناموفق از موضوعات بسیار داغ در حوزه هوش مصنوعی بوده که مکانیزم مشهور یادگیری با نگاه به گذشته (\lr{HER}) \cite{andrychowicz2017hindsight} و مدل متأخر با رویکرد پاداش نرم‌تر (\lr{Soft HER}) \cite{he2020softher} با استفاده از استراتژی برچسب‌گذاری مجدد هدف برای تجربیات ناموفق توانسته یادگیری محیط‌ها را سریع‌تر نمایند.

\textbf{نقد رویکرد انطباق با این حوزه:} برتری این مدل‌های عمیق کاملاً محدود به طراحی‌های وابسته به هدف و صریح (\lr{Goal-Conditioned Environments}) می‌باشد. به عبارت دیگر عامل باید یک مکان یا وضعیت شفاف (مانند موقعیت سه‌بعدی ابزار مکانیکی) را تصرف کند؛ در حالی که کنترل‌کننده‌های سیستم‌های خودوفقی اساساً و غالباً اهدافی به‌عنوان متغیر قطعی را دنبال نمی‌کنند، بلکه نگه‌داشتن انرژی یا تاخیر شبکه در حداقل‌ها یک سیاست آستانه‌محور است که با استراتژی تغییر برچسب و ارزیابی در بازتولید \lr{HER} با سازگاری اساسی تناقض دارد.

\section{سنتز و شناسایی شکاف تحقیقاتی}

با استناد به مطالب شرح داده شده و مروری جامع در کاستی مقالات پیشین، درمی‌یابیم که هیچ‌یک از دسته‌بندی‌های فعلی قادر به برطرف ساختن کل چالش‌ها به صورت جامع نبوده‌اند: راهکارهای یادگیری ماشین به پدیده‌های پیش‌بینی‌نشده محیط به دلیل فقدان آموزش آفلاین به شدت حساس و شکننده‌اند؛ مسیرهای اکتشاف جستجو با تاخیر غیرقابل پذیرش روبرو می‌باشند. حتی اگر از رویکرد‌های \lr{DRL} پایه بهره بگیریم، مستعد توقیف در استثنائات محلی بوده و قابلیت فراگیری جامع در مقابل نوسانات متغیر محیطی را متجلی نمی‌کنند. علاوه بر این ابزارهای پیشرفته مانند تخصیص مقاصد بازپخش گذشته امکان استفاده‌ی کارآمد مبتنی بر اهداف نامشخص سیستمِ خودوفقی را دارا نیستند.

بر این اساس یک \textbf{شکاف علمی‌و عملی مشخص} برای ابداع چارچوبی مبرهن مطرح است که دارای ویژگی‌های زیر باشد:
۱) تسلط اثربخش در فضاهای وسیع و گسسته وفقی بدون استفاده از جداول سنتی مقیاس ناپذیر،
۲) فراگیری هوشمندانه برخط بدون اجبار استفاده از داده‌های آموزش دیده‌ی گذشته،
۳) مقابله سرسختانه با بهینه‌های محلی بدون از دست دادن همگرایی و افت در محیط تغییریافته (نیاز به شکل‌دهی استراتژی آموزش)، و
۴) قابلیت پیوستگی با نیازهای غیرخطی یک سیستم خودوفقی بدون تغییر پارادایم تحمیلی به مسائلی متمرکز به تعیین هدف.

پیشنهاد بنیادی ما در این رساله مبتنی بر ایجاد یک مکانیزم جدید تلفیقی تحت عنوان \textbf{یادگیری تقویتی عمیق با شکل‌دهی تصادفی پاداش (\lr{RS-DRL})} است تا خلأ یاد شده کاملاً پر گشته و گامی‌کاربردی جهت ارتقای استقلال عوامل کنترلی نرم‌افزارهای پیچیده میسر گردد.

\section{خلاصه فصل}

این فصل در آغاز، با ترسیم بنیان‌های نظری مفاهیم یک سیستم پایدار و مکانیزم \lr{MAPE-K} و تشریح ماهیت مدل‌سازیِ عدم‌قطعیت با ساختار زمان‌اجرا بنا گردید. به تبع آن ضمن تشریح کارکردهای یادگیری تقویتی متدهای عمیق پیوستگی استراتژی پیشنهادی، محدوده‌های فعلی در مقابله با چالش مقیاس‌پذیری نمایان شد. در خاتمه، بر پایه ارزیابیِ تحلیلی در سابقه کارهای تحقیقاتی پیشین مشخص گردید، فقدان یک مکانیزم هوشمند برای شکل‌دهی تصادفی برخط کاملا مشهود است که در فصل آتی با مدل سیستم و فرمول‌سازی مسئله پیش‌نیازهای ساخت آن تنظیم می‌گردد.

\newpage

% فصل ۳: مدل سیستم و فرمول‌بندی مسئله
\chapter{مدل سیستم و فرمول‌بندی مسئله}
\section{مقدمه}

برای توسعه یک راهکار مبتنی بر یادگیری تقویتی به منظور مدیریت عدم قطعیت در سیستم‌های خودوفقی، ابتدا باید یک مدل صوری از مسئله وفقی‌سازی (تطابق) بنا نهیم. این فصل، مدل سیستم و فرمول‌بندی ریاضیاتی که بستر رویکرد پیشنهادی ما (\lr{RS-DRL}) را تشکیل می‌دهد، ارائه می‌نماید. ما با معرفی \lr{DeltaIoT} -- یک نمونه بارز از شبکه‌های حسگر بی‌سیم که چالش‌های کلیدی فضاهای وفقی بزرگ و عدم قطعیت‌های محیطی را به صورت آشکار به نمایش می‌گذارد -- بحث را آغاز می‌کنیم.

سپس مسئله وفقی‌سازی را به‌عنوان یک فرآیند تصمیم‌گیری مارکوف (\lr{MDP}) فرمول‌بندی کرده و فضای حالت، فضای عمل و تابع پاداش را برای سناریوهای بهینه‌سازی تک‌هدفه و چندهدفه تعریف می‌کنیم. نهایتاً، مسئله یادگیری یک سیاست وفقی بهینه را که می‌بایست در شرایط متغیر زمان اجرا و به‌طور مداوم عمل کند، به صورت صوری بیان می‌نماییم.

\section{مطالعه موردی انگیزشی: سیستم \lr{DeltaIoT}}

\subsection{مرور کلی سیستم و منابع عدم قطعیت}

سیستم \lr{DeltaIoT} یک شبکه حسگر بی‌سیم (\lr{WSN}) با تعاملات پیچیده است که به صورت گسترده به‌عنوان یک محک استاندارد (\lr{Benchmark}) برای ارزیابی سیستم‌های خودوفقی به کار می‌رود \cite{RS_DRL_Article}. این شبکه بسته به نسخه مورد استفاده، مابین ۱۵ تا ۳۷ گره (سنسور) دارد. دو منبع اصلی عدم قطعیت در این سیستم در زمان اجرا بروز می‌یابند:
\begin{enumerate}
    \item \textbf{نویز و تداخل شبکه (\lr{Network Interference/Noise}):} تغییرات محیطی و مداخلات پیرامونی که موجب نوسانات پیش‌بینی‌نشده در کیفیت ارتباطات از \lr{-40dB} تا \lr{+15dB} می‌گردد.
    \item \textbf{تغییرپذیری بار پیام‌ها (\lr{Message Load Variability}):} نوسان در تعداد پیام‌های تولید شده توسط هر سنسور (بین ۰ تا ۱۰ پیام) با الگوهای انتقال متفاوت در طول زمان.
\end{enumerate}

همانطور که در مطالعات پیشین (و روند نشان داده شده در شکل ۲ مقاله مرجع) نمایان است، این عدم قطعیت‌ها باعث شکل‌گیری افت شدید عملکرد سیستم در طی چرخه‌های وفقی گشته و اهمیت حیاتی مکانیسم‌های تطابق پویا را برجسته می‌کنند.

\subsection{فضای وفقی و بازنمایی پیکربندی}

فضای وفقی (\lr{Adaptation Space}) دربرگیرنده مجموعه تمامی‌گزینه‌ها و پیکربندی‌های ممکنی است که سیستم می‌تواند در زمان اجرا اتخاذ نماید. در \lr{DeltaIoT}، سه نسخه متفاوت با مشخصاتی به شرح زیر طراحی گردیده است \cite{RS_DRL_Article}:
\begin{itemize}
    \item \textbf{\lr{DeltaIoTv1}:} شامل ۱۵ سنسور و دارای ۲۱۶ گزینه وفقی متمایز.
    \item \textbf{\lr{DeltaIoTv1.1}:} شامل ۱۶ سنسور و دربرگیرنده ۱۲۹۶ گزینه وفقی است.
    \item \textbf{\lr{DeltaIoTv2}:} نسخه وسیع‌تر این شبکه متشکل از ۳۷ سنسور که دارای فضای عظیم متشکل از ۴۰۹۶ تصمیم ممکن می‌باشد.
\end{itemize}

هریک از این گزینه‌های وفقی در حقیقت ترکیبی از پارامترهای مختلف هستند؛ نظیر سطوح توان انتقال انرژی (\lr{Transmission Power}) برای سنسورهای مشخص و توزیع نسبی مسیرهای شبکه‌ای (\lr{Link Distribution}). انتخاب یک ترکیب از این پارامترها، یک پیکربندی کاملاً جدید و واحد را شکل می‌دهد.

\section{مدل صوری برای تصمیم‌گیری خودوفقی تحت عدم قطعیت}

برای امکان‌پذیر کردن یادگیری تقویتی، نیاز به فرمول‌بندی مسأله تطابق شبکه در قالب یک فرآیند تصمیم‌گیری مارکوف (\lr{MDP}) داریم که نمایانگر چهارگانه $(S, A, R, P)$ است \cite{RS_DRL_Article}. 

\subsection{تعریف فضای حالت}

فضای حالت مدل، وضعیت فعلی و پویای سیستم را با اندازه‌گیری سه معیار حیاتی عملکرد بازتاب می‌دهد. به بیان رسمی:
$$S = \{(E, P, L) \mid E \in \mathbb{R}, P \in [0,100], L \in [0,100]\}$$
که در آن:
\begin{itemize}
    \item \textbf{انرژی مصرفی ($E$):} مقدار حقیقیِ نمایانگر انرژی کل مصرفی توسط تمام سنسورهای شبکه در چرخه پیشین است.
    \item \textbf{تلفات بسته‌های شبکه ($P$):} درصدی از پیام‌های تولید شده که در طول مسیر به مقصد (دروازه اصلی) از دست رفته‌اند (بین ۰ تا ۱۰۰ درصد).
    \item \textbf{تاخیر ارسال ($L$):} معیاری معادل زمان مورد نیاز جهت رسیدن اطلاعات از سنسورها به گیت‌وی (حداکثر بازه ۱۰۰).
\end{itemize}

این سه متغیر مقداری به این دلیل انتخاب شده‌اند که تصویر کاملی از نیازمندی‌های کیفی (\lr{QoS}) شبکه را در یک سیستم توزیع‌شده با منابع محدود ارائه کرده و ارزیابی صحیحی برای عامل تصمیم‌گیرنده (\lr{Agent}) فراهم می‌کنند.

\subsection{فضای عمل: گزینه‌های وفقی گسسته}

فضای عمل شامل تمامی‌استراتژی‌ها و پیکربندی‌های ممکن به صورت مجموعه‌ای از گزینه‌های گسسته تعریف می‌شود:
$$A = \{a_1, a_2, \ldots, a_n\}$$
همان‌طور که پیش‌تر اشاره شد، مقدار $n$ بسته به ابعاد شبکه \lr{DeltaIoT} برابر با ۲۱۶، ۱۲۹۶ یا ۴۰۹۶ در نظر گرفته می‌شود. هر عمل $a_i$ نمایانگر یک پیکربندی مشخص از پارامترهای $\{p_{i1}, p_{i2}, \ldots, p_{im}\}$ است که در آن هر پارامتر $p_{ij}$ می‌تواند از مجموعه مقادیر محدود $\{v_{ij1}, v_{ij2}, \ldots, v_{ijk}\}$ تبعیت کند. این پارامترها در واقع توان ارسال سنسورها، توزیع لینک‌ها و مسیر یابی هر سنسور را تغییر می‌دهند.

\subsection{فرمول‌بندی تابع پاداش}

محاسبه صحیح تابع پاداش برای هدایت سیستم به سمت مقاصد از پیش تعیین شده، عاملی کلیدی محسوب می‌شود. بر اساس ماهیت سیستم و اهداف، ممکن است با بهینه‌سازی‌های تک‌هدفه، چندهدفه و یا اهداف آستانه‌ای مواجه شویم.

\subsubsection{فرمول‌بندی پاداش چندهدفه}

برای متوازن‌سازی رفتار سیستم حین برخورد با اهداف متعارض در طول چرخه خودوفقی، یک تابع پاداش بهینه‌سازی چندهدفه ارائه می‌گردد:
$$R = w_E \cdot \left(\frac{E_{max} - E}{E_{max} - E_{min}}\right) + w_P \cdot \left(\frac{P_{max} - P}{P_{max} - P_{min}}\right) + w_L \cdot \left(\frac{L_{max} - L}{L_{max} - L_{min}}\right)$$
در این نگاشت، ضرایب وزن‌دهی ($w_E + w_P + w_L = 1$) میزان اهمیت هر هدف را نشان می‌دهند و هر ترم موجود با استفاده از کران‌های بالا و پایین هر معیار محدودیت‌ها را بین بازه $[0, 1]$ نرمال‌سازی می‌کند.

\subsubsection{فرمول‌بندی پاداش تک‌هدفه}

زمانی که اولویت طراحی تنها برای غلبه بر یک مشکل واحد تعریف گردیده (مانند بهینه‌سازی صرفاً مصرف انرژی)، تابع فوق بدین شکل ساده‌سازی می‌گردد:
$$R = \left(\frac{E_{max} - E}{E_{max} - E_{min}}\right) \cdot w_E$$
در این معادله بقیه وزن‌ها معادل صفر در نظر گرفته شده‌اند که سبب هدایت بدون واسطه سیستم برای ارتقای سریع ویژگی هدف‌گذاری شده می‌شود.

\subsubsection{پارامترهای پاداش برای اهداف آستانه‌ای (\lr{TTS} و \lr{TT})}

برای مسائلی که نیازمند کنترل مطلق اهدافی در سطح استثناهای کیفیِ قابل قبول (\lr{Thresholds}) هستند (که در جدول ۲ مقاله با نام‌های \lr{TTS} و \lr{TT} نام برده شده‌اند)، تخصیص مقادیر پاداشِ مشروطی به صورت دودویی (\lr{Binary}) اعمال می‌گردد:
$$R = \begin{cases} 
1 & \text{اگر } P < \theta_P \text{ و } L < \theta_L \\
0 & \text{در غیر این صورت}
\end{cases}$$
که در آن مقادیر مرزی $\theta_P$ و $\theta_L$ به‌عنوان آستانه‌های از پیش تعیین شده معرفی شده‌اند. جدول \ref{tab:deltaiot_goals} خلاصه مقادیر \lr{TTS} و \lr{TT} (تنظیمات اهداف آستانه‌ای و ترکیبی) را پیرامون معماری‌ این نسخه‌ها نمایش می‌دهد:

\begin{table}[h!]
\centering
\caption{اهداف کیفی و آستانه‌ای (\lr{TTS} و \lr{TT}) سیستم \lr{DeltaIoT} \cite{RS_DRL_Article}}
\label{tab:deltaiot_goals}
\begin{tabular}{|c|c|c|c|c|}
\hline
\textbf{نسخه \lr{DeltaIoT}} & \textbf{تنظیمات هدف} & \textbf{تلفات بسته (\lr{Packet Loss})} & \textbf{تاخیر (\lr{Latency})} & \textbf{انرژی (\lr{Energy})} \\ \hline
\lr{DeltaIoTv1} & \lr{TTS} & کمتر از ۱۰٪ ($<10\%$) & کمتر از ۵٪ ($<5\%$) & کمینه‌سازی \\ \hline
\lr{DeltaIoTv1} & \lr{TT} & کمتر از ۱۰٪ ($<10\%$) & - & کمینه‌سازی \\ \hline
\lr{DeltaIoTv2} & \lr{TTS} & کمتر از ۱۰٪ ($<10\%$) & کمتر از ۵٪ ($<5\%$) & کمینه‌سازی \\ \hline
\end{tabular}
\end{table}

\section{بیان مسئله: یادگیری سیاست وفقی بهینه با عملیات مداوم در زمان اجرا}

با در نظر داشتن فرمول‌بندی \lr{MDP} در نظر گرفته شده، مسأله نهایی ما یادگیری روال سیاست وفقی بهینه $\pi^*: S \rightarrow A$ است. انتظار سیستم از این سیاست به حداکثر رسانی ارزش پاداش‌های تجمعی تنزیل‌یافته (\lr{Cumulative Discounted Reward}) است:
$$\pi^* = \arg\max_{\pi} \mathbb{E}\left[\sum_{t=0}^{\infty} \gamma^t R(s_t, a_t) \mid \pi\right]$$
در معادله فوق، مقدار $\gamma \in [0,1)$ ضریب تنزیل را معرفی می‌نماید. عملکرد موثر این سیاست بهینه باید همواره به‌طور مداوم و لاینقطع در زمان اجرا عمل کند و مکرراً با چالش‌های زیر روبرو گردد:
\begin{enumerate}
    \item \textbf{عدم قطعیت محیطی:} بار متغیر پیام‌ها و همچنین تداخل‌های ناپایدار و اختلال متناوب خطوط.
    \item \textbf{فضای عمل گسسته گسترده:} مدیریت فضای مقیاس‌پذیری که بسته به نسخه شبکه می‌تواند تا سقف ۴۰۹۶ تصمیم ممکن وسیع باشد.
    \item \textbf{پویایی غیرایستا (\lr{Non-stationary dynamics}):} نیاز به تطبیق مداوم و عدم کاهش سرعت واکنش به علت تغییرات شرایط از یک چرخه وفقی به چرخه بعدی (تخریب و کاهش عملکرد نشان داده شده در شکل ۲ مقاله مرجع).
\end{enumerate}

\section{خلاصه فصل}

فرآیند صورت‌بندی و مدل‌سازی در این فصل، زمینه‌ساز توسعه راهکار عملی این رساله گشته است. ما در این بخش ابتدا مطالعه موردی سیستم \lr{DeltaIoT} را شرح داده و ابعاد گوناگون فضای وفقی گسسته‌ی آن را تشریح نمودیم. در ادامه این مسأله با تمرکز به مدل تصمیم‌گیری مارکوف ساختاریافته شامل فضای وضعیت شبکه، اعمال تصمیم‌سازی و توابع پاداش‌های چندگانه برای بهینه‌سازی فرموله گردید.

تکمیل این چارچوب رسمی، قابلیت و ظرفیت یادگیری استراتژی‌های وفقی و منطبق با متد یادگیری تقویتی را فراهم می‌آورد. این پایه‌گذاری با اهمیت به ویژه به واسطه توصیف اهداف مشخص، در فصل آتی (به‌عنوان طراحی متد ابتکاری \lr{RS-DRL}) زمینه‌ساز توسعه راهکارِ خروج از بهینه‌های محلی با بهره‌گیری از شکل‌دهی تصادفیِ پاداش برای مدیریت عدم قطعیت خواهد بود.

\newpage

% فصل ۴: روش پیشنهادی \lr{RS-DRL}: معماری و مدل فرآیندی
\chapter{روش پیشنهادی \lr{RS-DRL}: معماری و مدل فرآیندی}
\section{مقدمه}

در فصل پیشین نشان دادیم که مدل‌های مرسوم یادگیری تقویتی، توانایی لازم برای پاسخ‌گویی به فضاهای بسیار بزرگ و گسسته تطابق را در زمان اجرای سیستم‌های خودوفقی ندارند و رویکردهای موجود اغلب در قله‌های محلی (\lr{Local Optima}) متوقف می‌شوند. همچنین آموزش مجدد مدل در زمان اجرا فرآیندی پرهزینه است و اتلاف وقت قابل توجهی برای یافتن گزینه‌های بهینه می‌طلبد. برای غلبه بر این محدودیت‌ها، سیستم نیازمند رویکردی است که بتواند با کمترین میزان نیاز به داده‌های از پیش‌تعیین‌شده، تجربیات پیشین خود را ارزیابی و استراتژی‌های جدیدی را در زمان اجرا به کار گیرد.

این فصل معماری و مدل فرآیندی سیستم خودوفقی مبتنی بر \lr{RS-DRL} را ارائه می‌دهد و تمرکز آن بر سازماندهی سطح سیستم، جریان کنترل زمان اجرا و الگوهای تعاملی مؤلفه‌ها است. در این راستا، ما از یک ساختار فرآیندگرا محور استفاده می‌کنیم—هر بخش پاسخ می‌دهد: «در مرحله بعد چه اتفاقی در سیستم می‌افتد و چرا؟» همچنین حلقه \lr{MAPE-K} به‌عنوان یک لنز توضیحی عمل می‌کند. این رویکرد به جای ارائه یک مدل یکپارچه پیچیده، از تصاویر تجزیه شده (\lr{Decomposed Figures}) برای توضیح روشن هر مرحله از فرآیند ۹ مرحله‌ای چرخه تطبیق استفاده می‌نماید.

\section{معماری سطح سیستم و حلقه تطبیق (\lr{Adaptation Loop})}

\subsection{انتزاع سیستم مدیریت‌شده و سیستم مدیریت‌کننده}

سیستم خودوفقی مبتنی بر \lr{RS-DRL} از دو جز کلیدی تشکیل می‌شود: \textbf{سیستم مدیریت‌شونده (\lr{Managed System})} (که منطق کاربردی را اجرا می‌کند، در اینجا شبکه \lr{DeltaIoT}) و \textbf{سیستم مدیریت‌کننده (\lr{Managing System})} (که منطق تطبیق را با استفاده از \lr{MAPE-K} تقویت‌شده با \lr{RS-DRL} پیاده‌سازی می‌نماید).

\subsection{چرخه پیوسته تطبیق (مراحل ۱ تا ۹)}

سیستم از طریق یک جریان تطبیقی ۹ مرحله‌ای عمل می‌کند:
\begin{enumerate}
    \item \textbf{مرحله ۱:} جمع‌آوری داده‌های زمان اجرا (\lr{Runtime data collection})
    \item \textbf{مرحله ۲:} ارسال وضعیت به ماژول \lr{RS-DRL}
    \item \textbf{مرحله ۳:} آموزش ناهمگام مدل و ذخیره‌سازی
    \item \textbf{مرحله ۴:} ارسال وضعیت به تحلیل‌گر (\lr{Analyzer})
    \item \textbf{مرحله ۵:} راه‌اندازی تطبیق (\lr{Adaptation triggering})
    \item \textbf{مرحله ۶:} بازیابی مدل آموزش‌دیده
    \item \textbf{مرحله ۷:} پیش‌بینی گزینه تطبیقی
    \item \textbf{مرحله ۸:} بازیابی جزئیات گزینه تطبیقی
    \item \textbf{مرحله ۹:} اعمال تطبیق (\lr{Adaptation application})
\end{enumerate}

همانطور که در شکل \ref{fig:rs-drl-architecture} مشاهده می‌شود، این معماری نشان‌دهنده یکپارچه‌سازی \lr{RS-DRL} با چارچوب کنترل \lr{MAPE-K} است.

\begin{figure}[H]
    \centering
    \adjustbox{width=0.95\textwidth,height=0.85\textheight,center,keepaspectratio}{\includegraphics{image/RS-DRL.png}}
    \caption{معماری یکپارچه‌سازی \lr{RS-DRL} با چارچوب \lr{MAPE-K}}
    \label{fig:rs-drl-architecture}
\end{figure}

\subsection{جداسازی تطبیق زمان اجرا و یادگیری ناهمگام}

یک تصمیم مهم طراحی معماری، جداسازی زمانی بین تطبیق زمان اجرا (مراحل ۱، ۴-۹) و یادگیری ناهمگام (مراحل ۲-۳) است. این جداسازی تضمین‌کننده پایداری در زمان اجراست، در حالی که امکان بهبود پیوسته سیاست‌ها را فراهم می‌آورد.

\section{نگاشت به مدل مرجع \lr{MAPE-K}}

\subsection{نگاشت مفهومی}

این بخش نگاشت مفهومی‌معماری \lr{RS-DRL} به مدل \lr{MAPE-K} را فراهم می‌آورد. این نگاشت به‌عنوان یک لنز طبقه‌بندی عمل می‌کند.
\begin{itemize}
    \item \textbf{پایش (\lr{Monitor}):} (مراحل ۱، ۲، ۴) مشاهده سیستم مدیریت‌شده و ساخت بازنمایی از فضای حالت.
    \item \textbf{تحلیل (\lr{Analyze}):} (مراحل ۴، ۵) ارزیابی وضعیت در برابر اهداف و راه‌اندازی تطبیق.
    \item \textbf{برنامه‌ریزی (\lr{Plan}):} (مراحل ۵، ۶، ۷) انتخاب گزینه‌های تطبیقی با استفاده از سیاست‌های \lr{RS-DRL}.
    \item \textbf{اجرا (\lr{Execute}):} (مراحل ۸، ۹) اعمال تطبیق‌های انتخاب‌شده.
    \item \textbf{دانش (\lr{Knowledge}):} (مراحل ۳، ۶، ۸) ذخیره‌سازی مدل‌ها، گزینه‌ها و تاریخچه وضعیت.
\end{itemize}

\subsection{نقش \lr{RS-DRL} در \lr{MAPE-K}}

مدل \lr{RS-DRL} به‌عنوان یک ارتقا در فاز \textbf{برنامه‌ریزی (\lr{Plan})} قرار می‌گیرد و هوش یادگیری را برای انتخاب گزینه تطبیقی فراهم می‌کند. این ماژول یک سرویس تصمیم‌گیری داخلی درون برنامه‌ریز است و نه کل معماری سیستم.

\section{پایش زمان اجرا و ساخت فضای حالت}

\subsection{جمع‌آوری داده‌ها، پیش‌پردازش و استخراج ویژگی}

فرآیند پایش با جمع‌آوری داده‌های زمان اجرا از سیستم مدیریت‌شده و محیط (مرحله ۱) آغاز می‌شود. داده‌های خام جمع‌آوری شده شامل خوانش سنسورها (مانند میزان انرژی، وضعیت لینک‌ها) می‌باشد. پس از جمع‌آوری، داده‌های خام پیش‌پردازش شده و ویژگی‌های مرتبط با اهداف استخراج می‌شوند.

\subsection{ساخت بردار حالت و اعتبارسنجی}

همانطور که در فصل قبل بحث شد، فضای حالت صوری شبکه شامل مفاهیمی‌پیرامون مصرف انرژی، اتلاف بسته‌ها و زمان تأخیر است:
$$S = \{(E, P, L) \mid E \in \mathbb{R}, P \in [0,100], L \in [0,100]\}$$
ویژگی‌های استخراج‌شده در این بردار حالت $s_t$ تجمیع می‌شوند که به‌عنوان رابط بین پایش و تصمیم‌گیری عمل می‌کند. این فرآیند با ارسال وضعیت معتبر به پایان می‌رسد.

\section{تحلیل زمان اجرا و راه‌اندازی تطبیق}

\subsection{ارزیابی اهداف و راه‌اندازی}

تحلیل‌گر (\lr{Analyzer}) بردارهای حالت را از مؤلفه پایش دریافت می‌کند (مرحله ۴) و ارزیابی می‌نماید که آیا وضعیت فعلی سیستم با اهداف تطبیقی سازگار است یا خیر. در صورتی که انحرافی از اهداف تضمین‌شده مشاهده شود یا اهداف در خطر نقض باشند، تحلیل‌گر فاز برنامه‌ریزی را راه‌اندازی می‌نماید (مرحله ۵).

\section{برنامه‌ریزی زمان اجرا و پیش‌بینی گزینه تطبیقی}

\subsection{بازیابی مدل و ارزیابی گزینه‌ها}

پس از دریافت تریگر تطبیق (مرحله ۵)، فرآیند برنامه‌ریزی، مدل آموزش‌دیده \lr{RS-DRL} را از مخزن دانش (مرحله ۶) بازیابی می‌کند.
سپس، پیش‌بینی‌کننده گزینه تطبیقی (\lr{Adaptation Option Predictor})، تمام گزینه‌های تطبیقی موجود (فضای عمل گسسته نظیر ۲۱۶، ۱۲۹۶، یا ۴۰۹۶ گزینه) را با استفاده از مدل بازیابی‌شده و مقادیر Q پیش‌بینی می‌نماید (مرحله ۷).
شایان ذکر است که این فرآیند در زمان اجرا یک \textbf{استنتاج خالص (\lr{Pure Inference})} است و هیچ‌گونه یادگیری یا نمونه‌برداری از حافظه در این مرحله انجام نمی‌شود که سبب تسریع در انتخاب گزینه بهینه می‌گردد.

\section{تایید گزینه تطبیقی (\lr{Adaptation Option Verification})}

\subsection{پارامترسازی مدل زمان اجرا و بررسی مدل آماری}

گزینه تطبیقی پیش‌بینی‌شده ابتدا نیاز به تایید دارد. این صحت‌سنجی با استفاده از مدل‌های زمان اجرای پارامتربندی‌شده با وضعیت فعلی و گزینه کاندید انجام می‌شود. ما از موتور \lr{UPPAAL-SMC} به همراه \lr{ActivFORMS} جهت بررسی آماری مدل (\lr{Statistical Model Checking}) استفاده می‌کنیم تا احتمال ارضای اهداف را تحت شرایط عدم قطعیت قبل از اعمال واقعی ارزیابی کنیم.

\section{اجرا و پیکربندی مجدد سیستم}

\subsection{بازیابی و اعمال تطبیق}

مؤلفه مجری (\lr{Executor}) مشخصات دقیق و تنظیمات پایین‌سطح برای گزینه تطبیقی انتخاب‌شده را از پایگاه دانش بازیابی می‌نماید (مرحله ۸). سپس، مجری این مشخصات را به تغییرات سیستمی‌در شبکه ترجمه کرده و بر روی سیستم مدیریت‌شونده اعمال می‌نماید (مرحله ۹). پس از آن، حلقه با از سرگیری پایش بسته می‌شود.

\section{یادگیری ناهمگام \lr{RS-DRL} و به‌روزرسانی سیاست}

\subsection{مروری بر چرخه یادگیری}

یادگیری به صورت کاملاً ناهمگام نسبت به تطبیق زمان اجرا پیش می‌رود و شامل جمع‌آوری تجربه‌ها، ذخیره‌سازی در حافظه بازپخش (\lr{Experience Replay})، آموزش مدل و ذخیره مجدد آن در پایگاه دانش (مرحله ۳) است.

\subsection{مکانیزم نوین شکل‌دهی پاداش (\lr{Reward Shaping})}

متمایزترین عامل نوآوری رساله، تعبیه مکانیزم شکل‌دهی پاداش تصادفی در طول فاز بازپخش تجربه است. برخلاف سیستم‌های معمول، سیستم ما تجربیات ناموفق گذشته ($r_i=0$) را به وسیله یک \textbf{فاکتور شکل‌دهی $\rho$} شناسایی و به‌عنوان تجربیات موفق تصنعی ($r_i=1$) بازطراحی می‌کند. این امر باعث می‌شود که عامل به جای گیر افتادن در سیاست‌های محلی محتاطانه، به اکتشاف بیشتر بپردازد.

الگوریتم \ref{alg:reward_shaping_v2} فرآیند دقیق این تغییرات بر روی مینی‌بچ‌ها را نمایان می‌سازد.

\begin{algorithm}[H]
\caption{شکل‌دهی پاداش (\lr{Reward Shaping})}
\label{alg:reward_shaping_v2}
\begin{algorithmic}[1]
\renewcommand{\algorithmicrequire}{\textbf{ورودی:}}
\renewcommand{\algorithmicensure}{\textbf{خروجی:}}
\REQUIRE مینی‌بچ $B$، فاکتور ضریبی شکل‌دهی $\rho$
\STATE $N \leftarrow |B|$
\STATE $K \leftarrow \rho \times N$ \COMMENT{محاسبه تعداد تجربیاتی که باید شکل‌دهی شوند}
\STATE شناسایی تجربیات ناموفقِ پیشین به شرط داشتن مقدار $r_i = 0$
\IF{تعداد کل تجربیات ناموفق مینی‌بچ $\geq K$}
    \STATE انتخاب تصادفی متعادل از میان $K$ تجربه ناموفق
    \FOR{هر تراکنش انتخاب شده نظیر $(s_i, a_i, 0, s_{i+1})$}
        \STATE $r_i \leftarrow 1$ \COMMENT{شکل‌دهی یک شکست به موفقیت تصنعی}
    \ENDFOR
\ELSE
    \STATE تمام تجربیات ناموفق حاضر را با پاداش ثابت ۱ شکل‌دهی کن
\ENDIF
\RETURN جریان مینی‌بچ به‌روزشده $B$
\end{algorithmic}
\end{algorithm}

\subsection{الگوریتم جامع آموزش \lr{RS-DRL}}

الگوریتم \ref{alg:rs_drl_v2} رویه کامل یادگیری در \lr{RS-DRL} را بر بستر معماری یکپارچه توصیف می‌کند.

\begin{algorithm}[H]
\caption{الگوریتم یکپارچه یادگیری \lr{RS-DRL}}
\label{alg:rs_drl_v2}
\begin{algorithmic}[1]
\STATE راه‌اندازی شبکه Q اصلی و شبکه تثبیت‌شونده هدف با وزن‌های مشابه $\theta$
\FOR{هر اپیزود تعاملی}
    \STATE دریافت حالت اولیه محیط $s$
    \WHILE{حلقه اپیزود تمام نشده است}
        \STATE انتخاب اقدام انطباقی $a$ با استفاده از سیاست $\epsilon$-گریدی
        \STATE اجرای اقدام $a$ بر روی محیط، و دریافت پاداش $r$ و حالت نهایی $s'$ (تولید سیگنال‌های پاداش از تاییدکننده)
        \STATE ذخیره‌سازی این تراکنش در حافظه بازپخشِ $\mathcal{D}$
        \STATE اخذ نمونه‌گیریِ یک مینی‌بچ رندوم $B$ از $\mathcal{D}$
        \IF{شرایط لازم برای بروزرسانیِ شکل‌دهی مهیا باشد}
            \STATE $B \leftarrow \text{RewardShaping}(B, \rho)$ \COMMENT{فراخوانی منطق شکل‌دهی}
        \ENDIF
        \STATE برآورد هدف تقربی $Q$ و اعمال نزول گرادیان جهت آموزش مدل
        \STATE انتقال درصدی وزن‌های شبکه هدف و تنظیم $s \leftarrow s'$
    \ENDWHILE
\ENDFOR
\end{algorithmic}
\end{algorithm}

\section{مدیریت دانش و هماهنگی فرآیندها}

\subsection{دانش به‌عنوان حافظه مشترک}

مخزن دانش (\lr{Knowledge}) به‌عنوان حافظه سیستمیک عمل کرده و مدل‌های آموزش‌دیده، گزینه‌های تطبیقی و تاریخچه وضعیت‌ها را نگه می‌دارد. این فرآیند با استفاده از قراردادهای خواندن و نوشتن همگام، به مدیریت ناهمگامی‌فعالیت‌های پایش (مرحله ۱)، آموزش مدل (مرحله ۳)، برنامه‌ریزی (مرحله ۶) و اجرا (مرحله ۸) کمک شایانی می‌کند.

\section{خلاصه فصل}

این فصل از طریق یک رویکرد تجزیه‌شده فرآیندی ساختاریافته، معماری تمام‌عیاری را جهت استقرار هوش مصنوعی تصمیم‌ساز در بطن یک شبکه مدیریت خودوفقی تبیین نمود. با تجزیه فرآیند به ۹ مرحله کلیدی از جمله پایش، تحلیل، حفظ، پیش‌بینی، تایید و اجرا، نشان دادیم که چگونه یک سیستم پیچیده \lr{MAPE-K} با یادگیری \lr{DRL} هماهنگ می‌شود. علاوه بر این، جداسازی زمانی بین تصمیم‌گیری زمان اجرا (پیش‌بینی خالص) و یادگیری ناهمگام (با استفاده از مکانیزم بدیع شکل‌دهی پاداش) سبب می‌شود خطاهای مدل‌سازی کمتری رخ داده و مدل قابلیت فرار از اکسترمم‌های محلی را بدون تأخیر در زمان واقعی پیدا کند. در فصل بعدی، به ارزیابی عملی و تجربیات به دست آمده از پیاده‌سازی این رویکرد خواهیم پرداخت.

% فصل ۵: ارزیابی تجربی
\chapter{ارزیابی تجربی}
\section{مقدمه}

ارزیابی تجربی نقش پررنگی در تأیید اثربخشی راهکارهای پیشنهادی در سیستم‌های خودوفقی ایفا می‌کند. به دلیل آنکه معماری نوین \lr{RS-DRL} ادعا می‌کند که می‌تواند بدون نیاز به آموزش‌های پیشرفته و آفلاین، با مدیریت تجربیات ناموفق، بر پدیده بهینه‌های محلی غلبه کند، اعتبارسنجی تجربی این دستاوردها بسیار حیاتی است. در این فصل با بهره‌گیری از سناریوهای متعدد ارزیابی، قابلیت‌های راهکار پیشنهادی را در پلتفرم معتبر شبکه حسگرهای بی‌سیم \lr{DeltaIoT} مورد سنجش جامع قرار می‌دهیم. 

این فصل در راستای پاسخگویی به سوالات تحقیق مطرح شده در فصل اول سازاماندهی گردیده و هدف اساسی آن نمایش و اثبات برتری این معماری نسبت به راهکارهای پایه و مرسوم موجود است. ساختار فصل بدین گونه است که در ابتدا اهداف ارزیابی تبیین گشته، سپس تنظیمات، مجموعه داده‌ها و محیط پیاده‌سازی معرفی می‌شوند؛ در گام‌های بعدی ضمن معرفی دقیق معیارهای سنجش، دستاوردها و نتایج به دست آمده در سناریوهای تک‌هدفه و چندهدفه و همچنین تحلیل مقیاس‌پذیری و ارزیابی آماری ارائه خواهد شد.

\section{اهداف ارزیابی و سوالات تحقیق}

انجام این ارزیابی تجربی همسو با پرسش‌های اساسی مطرح شده در فصل پیشین، اهداف مشخصی را دنبال می‌کند که در جدول \ref{tab:eval_rq} نقشه‌برداری شده‌اند:

\begin{table}[h!]
\centering
\caption{هم‌راستایی اهداف ارزیابی با سوالات تحقیق}
\label{tab:eval_rq}
\begin{tabular}{|p{7cm}|p{7cm}|}
\hline
\textbf{سوال تحقیق (\lr{RQ})} & \textbf{رویکرد ارزیابی} \\ \hline
\textbf{سوال اول:} طراحی مکانیزم \lr{DRL} برای خروج از بهینه‌های محلی & مقایسه منحنی‌های همگرایی (اشکال ۵ و ۶ در مقاله مرجع) در برابر روش‌های پایه و ارزیابی عملکرد مجانبی. \\ \hline
\textbf{سوال دوم:} ادغام شکل‌دهی پاداش با بازپخش تجربه & انجام مطالعه ابلیشن (\lr{Ablation Study}) در برابر گونه‌های مختلف از جمله \lr{HER} و \lr{Soft HER}. \\ \hline
\textbf{سوال سوم:} مدیریت فضاهای بزرگ و گسسته تطابق & اجرای تست مقیاس‌پذیری در بسترهای \lr{DeltaIoTv1} خ(۲۱۶ گزینه)، \lr{v1.1} (۱۲۹۶ گزینه) و \lr{v2} (۴۰۹۶ گزینه). \\ \hline
\textbf{سوال چهارم:} مقایسه \lr{RS-DRL} با روش‌های پیشرفته موجود & تقابل با روش‌های استانداردی نظیر \lr{$\epsilon$-greedy DQN}، \lr{DLASER+} و \lr{ML4EAS} (\lr{Tables 5-6}). \\ \hline
\end{tabular}
\end{table}

\section{تنظیمات آزمایشگاهی}

به منظور فراهم نمودن قابلیت بازتولید (\lr{Reproducibility}) آزمایش‌ها، جزئیات پیکربندی مطالعه موردی به دقت بیان می‌گردد. 

\subsection{نمونه‌های مطالعه موردی}
ما از پلتفرم \lr{DeltaIoT} به‌عنوان بستر آزمایشگاهی بهره برده‌ایم. پیکربندی استقراریافته برای شبکه‌های مورد بحث، مشخصات فضای حالت‌ها و پیچیدگی‌ها در جدول \ref{tab:case_study_instances} خلاصه گردیده است:

\begin{table}[h!]
\centering
\caption{نمونه‌های مورد مطالعه سیستم \lr{DeltaIoT}}
\label{tab:case_study_instances}
\begin{tabular}{|c|c|c|c|}
\hline
\textbf{نمونه (نسخه)} & \textbf{تعداد سنسورها} & \textbf{اندازه فضای وفقی (گزینه‌ها)} & \textbf{گام‌های آموزش} \\ \hline
\lr{DeltaIoTv1} & ۱۵ & ۲۱۶ & ۱۰,۰۰۰ \\ \hline
\lr{DeltaIoTv1.1} & ۱۶ & ۱,۲۹۶ & ۶۴,۸۰۰ \\ \hline
\lr{DeltaIoTv2} & ۳۷ & ۴,۰۹۶ & متناسب با مقیاس شبکه \\ \hline
\end{tabular}
\end{table}

\subsection{مجموعه داده‌ها}
جهت شبیه‌سازی شرایط در زمان اجرا، از ردپاهای محیطی (\lr{Traces}) مرتبط با سیستم \lr{DeltaIoT} استفاده شده که شامل اطلاعات چرخه‌های تطابق می‌باشند \cite{RS_DRL_Article}. این مجموعه داده، الگوهای تداخلات شبکه‌ای متغیر مابین $\text{-40dB}$ تا $\text{+15dB}$ را ثبت کرده و همچنین نوسانات مرتبط با بار ترافیک پیام‌ها (بین ۰ تا ۱۰ پیام به ازای هر سنسور) را در طول زمان شامل می‌شود.

\subsection{پلتفرم پیاده‌سازی}
پلتفرم عملیاتی اجراشده در این تحقیق شامل مجموعه‌ای از پیشرفته‌ترین چارچوب‌های محاسباتی است. برای ایجاد محیط یادگیری تقویتی از کتابخانه استاندارد \lr{Gymnasium} بهره گرفته شد و پیاده‌سازی شبکه‌های \lr{DQN} بر روی چارچوب مشهور \lr{Stable-Baselines3} انجام شده است. در دل این پیاده‌سازی، مکانیزم بومی‌شکل‌دهی پاداش تصادفی به صورت مستقیم به درون حافظه بازپخش (\lr{Replay Buffer}) تزریق گشته است. به منظور اعتبارسنجی صوری و ایمنی گزینه‌های تطبیق‌یافته نیز موتور \lr{UPPAAL-SMC} و مؤلفه واسط \lr{ActivFORMS} نقش رابط‌های کنترلی را ایفا می‌کنند.

\section{روش‌های پایه و مقایسه}

جهت ارزیابی عملکرد و تفوق رویکرد \lr{RS-DRL}، آن را با گستره‌ای از متدهای موجود بررسی نموده‌ایم. مقایسه در سه دسته‌بندی هدف‌گذاری به شرح زیر صورت گرفت:

\begin{itemize}
    \item \textbf{مقاصد کمینه‌سازی (\lr{Minimization Goals}):} برای این منظور مدل در برابر عامل استانداردِ یادگیری پایه، نظیر \lr{$\epsilon$-greedy DQN} و همچنین در برابر ترکیبات \lr{RS-DRL} با استراتژی‌های مشهوری چون \lr{HER} (\lr{RS-DRL-with-HER}) و \lr{Soft HER} قرار گرفت. 
    \item \textbf{مقاصد ترکیبی اهداف آستانه‌ای محدود (\lr{TTS Goal}):} برای ارزیابی کیفی، الگوریتم‌های شناخته‌شده‌ای همچون \lr{DLASER+} و یک سیستم مرجع (\lr{Reference}) که بر پایه جستجوی کامل و جامع تمامی‌حالت‌ها (\lr{Exhaustive Search}) استوار بود، مورد ارزیابی قرار گرفت.
    \item \textbf{مقاصد آستانه‌ای مطلق (\lr{TT Goal}):} راهکار پیشنهادی در تقابل با مدل‌های یادگیری پیشرفته همچون \lr{ML4EAS} (\lr{Machine Learning for Evaluating Adaptation Space})، نسخه اولیه \lr{DLASER}، \lr{DLASER+} و روش مرجع سنجیده شد.
\end{itemize}

\textbf{نکته کلیدی در خصوص انطباق روش‌های \lr{HER}:} ذکر این نکته حائز اهمیت است که مدل‌های \lr{HER} و \lr{Soft HER} اساساً برای محیط‌های مبتنی بر هدف‌های صریح (\lr{Goal-Conditioned Environments}) طراحی و تنظیم گشته‌اند. با این وجود برای فراهم کردن بستر مقایسه علمی‌در محیط‌های نرم‌افزاری شبکه همانند \lr{DeltaIoT}، ما این راهکارها را متناسب‌سازی کرده و با در نظر گرفتن آستانه‌های تعریف شده اهداف کیفی (\lr{QoS}) به‌عنوان اهداف جایگزین (\lr{Proxy Goals})، فرآیند تغییر مرجع مجدد (\lr{Relabeling}) را بر اساس آنها شبیه‌سازی کردیم.

\section{معیارهای ارزیابی}

در راستای متدولوژی‌های سنجش استاندارد موجود، ارزیابی این رساله بر محور دو گروه دسته‌بندی استوار گشته است:

\subsection{معیارهای عملکرد یادگیری عامل هوشمند}
این معیارها برای بررسی روند یادگیری سیستم استفاده می‌شوند:
\begin{enumerate}
    \item \textbf{عملکرد مجانبی (\lr{Asymptotic Performance}):} این سطح، وضعیت نهایی‌ترین ارزش و کیفیتی است که عامل پس از دریافت آموزش کافی به آن دست می‌یابد. اندازه‌گیری آن به‌طور معمول بر پایه محاسبه معدل پاداش به دست آمده روی تعداد محدودی از گام‌های نهایی چرخه امکان‌پذیر است.
    \item \textbf{زمان تا آستانه (\lr{Time to Threshold}):} بیانگر تعداد گام‌های زمانی است که عامل نیازمند طی کردن است تا به یک مقدار پاداش از پیش تعیین‌شده همگرا گردد. به‌عنوان نمونه آستانه مورد قبول در جداول ارزیابی مقداری معادل ۰.۹ در نظر گرفته شده است.
    \item \textbf{عملکرد کل (\lr{Total Performance}):} تابعی برای ارزیابی جامع که از مجذور عملکرد مجانبی در طی زمان کم شده از جمع کل مقادیر پاداش اخذ شده به دست می‌آید. این فرمول به صورت زیر نمایش داده می‌شود:
    $$TP = R_{\infty}\cdot T - \sum_{t=1}^{T} r_t$$
    \textbf{تفسیر این معیار:} به دست آمدن مقادیر مثبت نشانگر آن است که منحنی یادگیری مدل همواره پایین‌تر از سطح آرمانی مجانبی بوده و میل به بهبود و رشد در طول زمان دارد. اما مقادیر منفی حاکی از اوج‌گیری اولیه و سپس افت و همگرایی به سطوح پایین‌تر عملکرد در زمان اجراست.
\end{enumerate}

\subsection{معیارهای عملکرد سیستم}
موارد زیر منعکس‌کننده خروجی نهایی سیستم در ساختار شبکه می‌باشند:
\begin{itemize}
    \item \textbf{تلفات بسته‌های ارسالی (\lr{Packet Loss}):} درصد پیغام‌هایی که با خطا مواجه شده و به مقصد نمی‌رسند.
    \item \textbf{تاخیر ارسال (\lr{Latency}):} مدت زمان صرف‌شده برای انتقال اطلاعات به گره مرکزی.
    \item \textbf{مصرف انرژی (\lr{Energy Consumption}):} مصرف نهایی انرژی در گره‌های سنسورها پس از اعمال وفقی‌سازی.
\end{itemize}

\section{نتایج و تحلیل‌ها}

این بخش به‌عنوان قلب ارزیابی رساله، نتایج استخراج شده از شبیه‌سازی‌ها را تدوین و تفکیک نموده است.

\subsection{منحنی‌های یادگیری و رفتار همگرایی}

روند پیشرفت یادگیری الگوریتم در مقایسه با مدل پایه‌ای به وضوح برتری مکانیزم پیشنهادی را تصدیق می‌کند.
در تحلیل رفتار شبکه با استفاده از \lr{DeltaIoTv1}، مشاهده می‌گردد که رویکرد \lr{RS-DRL} توانسته است ارزش مجانبی بسیار درخشان ۰.۹۹۰۵ را (با خطای انحراف معیار ۰.۰۲۶۸) کسب نموده و در مدت زمان نسبتاً کوتاه ۲۲۶۶ گام (با تحمل خطای ۱۵۰ پیش و پس) به آستانه مطلوب نفوذ کند. در نقطه‌ی مقابل، الگوریتم پایه‌ی حریصانه یعنی \lr{$\epsilon$-greedy DQN} نه تنها نتوانست از حد مجانبی ۰.۴۵۵۹ فراتر رود، بلکه هرگز موفق به فتح آستانه‌ی در نظر گرفته شده نیز نگردید. دلیل بنیادین این توفیق \lr{RS-DRL}، آن است که مکانیزم شکل‌دهی تصادفی، به جای اکتفای تنها بر آموزش پایه، با تخصیص پاداش به تجربیات منحصراً ناموفق اجازه می‌دهد این مدل نه تنها به ثبات بالاتری بهبود یابد، بلکه در بالاترین سطح مجانبی با ثبات مطلوبی مستقر بماند و از بهینه‌های محلی فرار کند.

در فضای بسیار پیچیده‌تر با حضور ۱۲۹۶ گزینه تطبیق‌پذیر در بسته‌ی \lr{DeltaIoTv1.1}، راهکار ما در ۱۹۲۹۶ گام به پاداش عالی ۰.۸۹۸۳ (±۰.۰۹۳۲) نفوذ پیدا کرد، در حالی که همه مدل‌های رقیب عملکرد کل ضعیف‌تری را ثبت کرده و هرگز به خط آستانه نزدیک نشدند.

\subsection{مقایسه عملکرد به ازای هر هدف}

به منظور درک دقیق جزئیات کارکرد عوامل کنترلی روی تک‌تک متغیرهای سیستم، به ارزیابی رویکرد پیشنهاد شده با انواع \lr{HER} پرداختیم:

\begin{itemize}
    \item \textbf{بهینه‌سازی چندهدفه (\lr{Multi-objective}):} متد ما توانست به پاداشی معادل ۰.۹۹۰۵ (در شبکه اول) و ۰.۸۹۸۳ (در شبکه دوم) دست یابد در حالی که بهترین انواع \lr{HER} منحصراً مقادیری مابین ۰.۴۹ الی ۰.۶۱ را ثبت نمودند. علت این افت فاحش در مدل‌های \lr{HER} را باید در ماهیت وابستگیِ این مدل‌ها به تعاریف و تخصیص اهداف صریح، در مقایسه با پویاییِ اهداف نامشخص اکوسیستمِ \lr{SAS} جستجو نمود.
    \item \textbf{مصرف انرژی:} شاخص‌های ارزیابی، مصرف انرژی \lr{RS-DRL} را با ثبت میانگین ۰.۹۷۵۶ (در \lr{v1}) و ۰.۹۳۰۱ (در \lr{v1.1}) بسیار متمایز از مدل‌های پایه (بازه‌‌ی ۰.۳۹ الی ۰.۵۸) نشان داد.
    \item \textbf{تلفات بسته‌ها (\lr{Packet Loss}):} دستیابی راهکار توسعه‌یافته در نگه‌داری امنیت شبکه‌ها به‌عنوان عاملی در مرز نزدیک به حالت بهینه؛ ۰.۹۸۵۰ در شبکه اولیه و تغییر آن به وضعیت ۰.۹۱۲۹ در شبکه‌ی پیچیده‌تر دوم.
    \item \textbf{تاخیر (\lr{Latency}):} شبکه مجهز به \lr{RS-DRL} موفق‌ترین سرعت رکوردگیری را با معدل ۰.۹۹۶۴ و ۰.۹۰۲۱ منعکس کرد.
\end{itemize}

\subsection{ارزیابی مقایسه‌ای در برابر تمام روش‌های پایه}

ارزیابی شبکه در تقابل با دیگر رقبا نظیر متدهای جستجوگر، بسیار حائز اهمیت است:
\begin{itemize}
    \item \textbf{نتایج هدف \lr{TTS}:} تحت شبکه اولیه، متد \lr{RS-DRL} از دست رفتگی بسته‌های حاوی پیغام را در مرز ۱۰.۵۷۸٪ متوقف نمود، در حالی که استفاده از ارزیابی جستجوی مرجع (\lr{Reference}) درصد ۱۴.۲۳٪ و همچنین \lr{DLASER+} بالغ بر ۱۲.۸۹٪ اتلاف داشته‌اند. تاخیر ارسال با نرخ ۰.۰۹۷ بسیار بهتر از رقیب اصلیِ آن یعنی \lr{DLASER+} (با ثبت ۰.۸۴۰) کارنامه موفقی از خود برجای گذاشت و متوازن با روش مرجع (۰.۱۰۰) بود. میزانِ مصرف انرژی (با مقدار ۱۲.۸۸۲) نیز نزدیک به رویکرد مرجع (۱۳.۱۰۰) و \lr{DLASER+} (۱۳.۰۲۰) محاسبه گردید.
    \item \textbf{نتایج هدف \lr{TT}:} راهکار ما در پلتفرم \lr{DeltaIoTv1}، به تلفات بسیار کم شبکه (۷.۱۵۳٪) در مقابل بهترین رقیبش یعنی \lr{DLASER+} (۷.۹۸۰٪) رسید؛ همچنین در پلتفرم عظیم \lr{v2}، موفق شد مقدار ۹.۳۳۰٪ را برای این فاکتور ثبت نماید در حالیکه بهترین سیستم هوش ماشینی فعلی (\lr{ML4EAS}) مقداری معادل ۱۰.۵۶۰٪ ارائه داد.
\end{itemize}

تنها تفاوت موجود، تخصیص تاخیر در نسخه دوم بود (که \lr{RS-DRL} تاخیر ۰.۲۵۱ را در مقایسه با ۰.۱۰۰ روش مرجع ثبت نمود)؛ که البته یک طراحی مبتنی بر خط‌مشی به جهت اولویت‌بندی انتقال شبکه و کاهش اتلاف به جای کاهش زمان تاخیر قلمداد می‌شود و نه یک شکست الگوریتمی.

\subsection{مطالعه ابلیشن: حساسیت به فاکتور شکل‌دهی $\rho$}

یکی از دستاوردهای علمی‌مهم این متد، تخصیص استراتژیک ضریب مفقود به تجربیات گذشته بود. در الگوریتم، مقدار کنترل‌کننده $\rho$ تعیین می‌کند تا چه تعداد از تجربیات ناکامِ ثبت‌شده در مینی‌بچ (ذخیره‌شده در حافظه‌ی سیستم)، با ضریب عدد یک (پاداش مثبت تصنعی) جایگزین شوند. 

در صورتی که تجربه‌های شکست خورده موجود در حافظه بیشتر از نرخ انتخابی ما باشد، کاملاً به صورت تصادفی تعدادی انتخاب گشته و مجدد شکل‌دهی می‌گردند و در غیر این صورت روی تمامی‌آنها شکل‌دهی پیاده می‌گردد.
\textbf{نکات تحلیلی پیرامون فاکتور:}
ایجاد یک توازن کلیدی حول فاکتور بسیار مهم است. اگر $\rho$ بیش از حد بالا باشد، عامل با نویز بالا مواجه خواهد شد. متقابلاً مقادیر بسیار پایین نیز به آن معناست که عامل استراتژی جسورانه‌ای برای خروج از بهینه‌های محلی نخواهد داشت و مشابه شبکه‌های پایه فریز می‌شود. این انتخابِ تصادفی در کنار توزیع منطقی $\rho$، از تناسب‌بیش‌ازحد (\lr{Overfitting}) حافظه نیز ممانعت به عمل می‌آورد.

\subsection{تحلیل مقیاس‌پذیری}

مقایسه نتایج مدل حاکی از پیروزی پایدار آن است. تغییر از یک اکوسیستم ساده متشکل از ۲۱۶ گزینه به مقیاس بسیار وسیع با ۴۰۹۶ تصمیم ممکن، نشان داد مدل پیشنهادی رساله به نحوی قدرتمند برتری‌های خود را همچنان حفظ می‌کند و تنها زمان همگراییِ یادگیری (جهش از ۲۲۶۶ به ۱۹۲۹۶ گام در نسخه \lr{v1.1}) به تبع پویایی‌های اضافه شده افزایش یافته است.

\section{اعتبارسنجی آماری}

برای استواری یافته‌های این رساله، نتایج حاصله توسط آزمون‌های دقیق آماری اعتبارسنجی شده‌اند. 
روش‌شناسی توسعه‌یافته در ارزیابی شامل حداقل ۳۰ دور اجرای کاملاً مستقل با بذرهای رندوم، انجام آزمون تی مستقل تک‌نمونه‌ای (\lr{One-sample t-tests}) با ضریب خطای آلفای برابر ۰.۰۵، و تخصیص فواصل اطمینان ৯۵٪ به روش بوت‌استرپ است.
فرمول ریاضی محاسبه آماری به این شکل تنظیم گردید:
$$t = \frac{\bar{x} - \mu_0}{\frac{s}{\sqrt{n}}}$$

\textbf{خلاصه اعتبارسنجی‌ها (بر اساس مقادیر $p$):}
از بین مجموعاً ۲۸ مقایسه‌ی متمایز و شاخص، تعداد ۲۵ مقایسه کیفیتِ الگوریتم توسعه‌یافته در یک شرایط آماری فوق‌العاده قرار داشتند (۸۹.۳٪ تست‌ها معنادار بودند). شگفتی زمانی کامل گشت که در پارامتر اتلاف بسته‌های شبکه (\lr{Packet Loss}) تمام ۱۰۰٪ از تست‌ها ارزش بهبود معنادار ارائه دادند؛ به‌طور نمونه، احتمال وقوع تصادفی برتری آن برای متد \lr{TTS} شبکه اول در مقابل مدل‌های \lr{DLASER+} و رویکرد مرجع کمتر از ۰.۰۰۱ برآورد گردید (\lr{$p < 0.001$}). مواردی جزئی نیز که نتایج قابل توجهی نداشتند (مانند آزمون تاخیر که مقدار \lr{$p=0.134$} در شبکه اولیه دریافت نمود یا مصرف انرژی که در حالت \lr{TTS} مقداری معادل \lr{$p=0.326$} ثبت کرد) عمدتاً به دلیل عملکرد متعادل با روش مرجع بوده و هیچگونه شواهدی مبنی بر تضعیف متدِ موردِ بررسی به دست نمی‌دهد.

\section{خلاصه فصل}

این فصل به نحو جامعی توانست ادعاهای نوآوری رساله پیرامون مکانیزم شکل‌دهی یادگیری را تایید نماید. با مرور نتایج مشخص گردید: ۱) متد ابداعی \lr{RS-DRL} با پایداری ویژه‌ای توانسته تمامی‌الگوریتم‌های پرآوازه و قدرتمند فعلی پیشین را در ابعاد بهبود عملکرد یادگیری کیفیت و کنترل سیستم شکست دهد. ۲) این برتری در حدود ۸۹.۳٪ مواقع توسط تست‌های اعتبارسنجی آماری تأیید گردیده‌اند (۲۵ مورد از ۲۸ مقایسه کلیدی). ۳) این متد ثابت نمود تغییر مقیاس گزینه‌های وفقی از ۲۱۶ به ۴۰۹۶ گزینه نمی‌تواند قابلیت‌های آن را مختل نماید. ۴) مکانیزم شکل‌دهی مجدد توانسته در مقایسه دقیق و عمیق، از وقوع توقف سیستم‌ها در تله‌های استثناهای محلی جلوگیری کند و نهایتاً ۵) برخلاف متدهای رایج هم‌نوع پیشرفته خود نظیر مدل‌های \lr{HER}، برای پویاییِ پایدار اساساً نیازمند تخصیص دانش سیستم‌ها با دامنه نامشخص (\lr{Domain-specific knowledge}) و یا اهداف صریح از پیش مقدر نبوده است.

\newpage

% فصل ۶: بحث، کارهای آینده و نتیجه‌گیری
\chapter{بحث، کارهای آینده و نتیجه‌گیری}

\section{بحث در مورد یافته‌های کلیدی}

این بخش با ترکیب نتایج اصلی و تشریح اهمیت آن‌ها، یافته‌های مستخرج از رساله را مورد ارزیابی عمیق قرار می‌دهد.

\subsection{نظریه‌پردازی مکانیزم: چرا \lr{RS-DRL} کار می‌کند؟}

همان‌طور که در نتایج ارزیابی تجربی (مندرج در فصل ۵ و با استناد به شکل‌های ۵ و ۶ در مقاله مرجع \cite{RS_DRL_Article}) نشان داده شد، یکی از مهم‌ترین نقاط قوت \lr{RS-DRL} توانایی آن در فرار از بهینه‌های محلی است که یک چالش جدی و متداول در سیستم‌های \lr{DRL} محسوب می‌شود. در محیط‌های پویا، پاداش متوسط مدل \lr{RS-DRL} نه تنها به‌طور پیوسته افزایش می‌یابد، بلکه در وضعیت پایداری نسبت به رویکرد پایه \lr{$\epsilon$-greedy DQN} - که با نوسانات شدیدی همراه است و نمی‌تواند از راه‌حل‌های غیربهینه رهایی یابد - در سطحی بسیار بالاتر تثبیت می‌گردد.

اثربخشی رویکرد \lr{RS-DRL} را می‌توان به سه مکانیزم به هم‌پیوسته زیر نسبت داد:

\begin{enumerate}
    \item \textbf{اکتشاف از طریق تزریق پاداش خوش‌بینانه:} با تغییر تدریجی و تصادفیِ تجربیات ناموفق گذشته ($r_i=0$) به تجربیات موفق ($r_i=1$)، عامل یادگیرنده سیگنال‌های تقویت‌کننده مثبتی را از اعمال و وضعیت‌هایی دریافت می‌کند که در شرایط عادی از اولویت خارج می‌شدند. این امر از همگرایی زودهنگام مدل به سیاست‌های غیربهینه که یکی از معضلاتِ بنیادین \lr{DQN} استاندارد است، ممانعت به عمل می‌آورد.
    \item \textbf{شکستن همبستگیِ زمانی در الگوهای شکست:} انتخاب تصادفی در فرآیند تغییر تجربیات مردود (که توسط پارامتر $\rho$ کنترل می‌گردد) متغیری را معرفی می‌کند که مانع از حفظ بیش‌حد (\lr{Overfitting}) عامل بر روی شرایط متوالی شکست می‌شود. این متغیر به ویژه در محیط‌های غیرایستایی (\lr{Non-stationary}) که در آن‌ها شکست دیروز ممکن است استراتژی بهینه امروز تلقی گردد، بسیار راهگشا خواهد بود.
    \item \textbf{یادگیری مستقل از دامنه (\lr{Domain-Agnostic}):} برخلاف روش‌های متکی به \lr{HER} که مستلزم تعاریف صریح و قطعی از اهداف هستند، شکل‌دهی تصادفیِ \lr{RS-DRL} بدون هیچ‌گونه نیاز به درک معناییِ اهدافِ دامنه عمل می‌کند. این موضوع، اجرایِ آن را ذاتاً برای سیستم‌هایی انعطاف‌پذیر می‌سازد که با اهداف تلویحی و آستانه‌محور - که در سیستم‌های خودوفقی بسیار شایع می‌باشند - سروکار دارند.
\end{enumerate}

\subsection{تحلیل تعاملات و مصالحه‌ها (به‌عنوان مثال، تاخیر در مقابل از دست رفتن بسته)}

در جداول مقایسه‌ای پیشین (برای نمونه جدول ۶ در توصیفات ارزیابی)، مشاهده گردید که برای تخصیص هدف \lr{TTS} در شبکه \lr{DeltaIoTv1}، مدل \lr{RS-DRL} به درستی به تاخیر متوازنی دست پیدا کرد (مقدار ۰.۰۹۷ در برابر ۰.۱۰۰ روش مرجع) و در مجموع عملکرد بسیار بهتری در مقایسه با \lr{DLASER+} نشان داد. با این وجود، برای شبکه \lr{DeltaIoTv2} تحت هدف \lr{TT} مقدار تاخیر بالاتری حاصل شد (۰.۲۵۱ در برابر ۰.۱۰۰ روش مرجع).

افزایش تاخیر در نسخه دوم شبکه برای مقاصد مربوط به \lr{TT} نشان‌دهنده یک مصالحه (\lr{Trade-off}) بسیار مهم در عملکرد عامل است: هنگامی‌که هدف کلیدیِ تعریف‌شده از سوی کاربر (مثلاً حفظ توانایی توان عملیاتی شبکه) با اولویت‌بندی به حداقل رساندن مشکل فقدانِ بسته توأم گردد، مدلِ هوشمند درمی‌یابد که باید به منظور جلوگیری از دست‌افتادگیِ پیام‌ها، به سمتِ مسیرهای قابل اعتمادتری متمایل شود که به صورت ذاتی کندتر هستند. در محیط‌های بهینه‌سازی چندهدفه، چنین مصالحه‌ای نه تنها یک نقطه ضعف و محدودیت به شمار نمی‌رود، بلکه از قضا یک رفتار برآمده و مطلوب از سوی عامل ارزیابی می‌گردد. چرا که در اینجا اولویت سیستم بر تضمینِ رساندن بسته‌ها بیش از حفظ سرعت بوده است و در کاربردهایی که صحت اطلاعات در آن‌ها اهمیت حیاتی دارد (مانند مانیتورینگ زیست‌محیطی یا سیستم‌های کنترل صنعتی) این تصمیم به‌عنوان عاقلانه‌ترین انتخاب مهندسی تلقی خواهد شد.

این رفتار کاملاً با اصول انتخاب بهینه در استراتژی‌های پرتو-اپتیمال (\lr{Pareto Optimality}) همخوانی دارد. هنگامی‌که مرزهای متضاد انتخاب‌ها بررسی می‌گردند، تنظیمات دارای تأخیرِ بی‌نهایت کم‌تر عموماً ارتباط معناداری با نرخ بالای اتلاف بسته‌ها – به دلیل استراتژی‌های تهاجمی‌ترِ انتقال داده – نشان می‌دهند. راهکار \lr{RS-DRL} یاد می‌گیرد این خطِ مصالحه را با ظرافت طی کند و تصمیم‌ها را منطبق بَر بستر نیاز و با تخصیص وزن‌های اختصاص‌یافته برای هر متغیر اخذ نماید.

\subsection{پیامدهای مهندسی سیستم‌های خودوفقی}

ثبات، پایداری و همگرایی سریع متد \lr{RS-DRL} عوامل حیاتی در محیط‌های پویایی هستند که مقیاس و شرایط آن‌ها به سرعت در حال جابجایی است. از طریق در هم تنیدگی مکانیزم شکل‌دهی پاداش، \lr{RS-DRL} تضمین می‌نماید که عامل تصمیم‌ساز حتی بدون نیاز به بازآموزی و تخصیص پیکربندی‌های جدید به تغییرات پاسخ می‌دهد، در حالی که در اکثر \lr{DRL}های استاندارد الزام به طی مجدد مسیر آموزش برای اداره کردن امور تغییر یافته، مشهود و گریز‌ناپذیر است. این خصوصیت، این مدل را برای سامانه‌های برخط مانند سیستم‌های ابری و شبکه‌های اشیاء که استقرار و بروزرسانی خودکار استراتژی در آن‌ها واجب است، بسیار متمایز می‌سازد.

به بیان دقیق‌تر، پیامدهای عملی این رهیافت برای مهندسین توسعه‌دهنده محیط‌های خودوفقی شامل موارد زیر است:

\begin{enumerate}
    \item \textbf{کاهش قابل‌ملاحظه سربار عملیاتی:} مدل‌های سنتی به هنگام برخورد با رانش و تغییر شرایط محیطی نیازمند صرف چرخه‌های متعدد بازآموزیِ پیاپی می‌باشند. در حالیکه ویژگی‌های تعمیم‌پذیری بالای \lr{RS-DRL} امکان دور ماندنِ مستمرِ مدل از بروزرسانی‌های مجدد را فراهم می‌سازد که خود علاوه بر صرفه‌جویی شدید محاسباتی باعث جلوگیری از عملکرد ضعیف مدل در طی ادوارِ گذارِ بازآموزی می‌گردد.
    \item \textbf{بی‌نیازی از داده‌های برچسب‌دار:} بر خلاف مدل‌هایی همانند \lr{ML4EAS} یا \lr{DLASER+} که اتکای گسترده‌ای به تجمیعِ پیشینِ داده‌ها داشته‌اند، الگوریتم \lr{RS-DRL} کاملاً آنلاین و درلحظه توسعه می‌یابد. این قابلیت نقطه قوتِ برجسته‌ای در محیط‌هایی محسوب می‌گردد که تهیه اطلاعات پیش‌آموزشی غیرممکن و پیش‌بینی‌ناپذیر است.
    \item \textbf{پایداری و کاهش تخریب تدریجیِ کیفیت متناسب با گسترش:} استواریِ معنادارِ کیفیت و عملکرد در مواجهه با نسخه‌ی ۲۰۱۶-گزینه‌ای (نسخه \lr{v1})، گسترش ۱۲۹۶-گزینه‌ای (نسخه \lr{v1.1}) و مقیاس سرسام‌آور ۴۰۹۶ حالت ممکن (نسخه \lr{v2}) مؤید آن است که این راهکار متناسب با میزان گستردگی شبکه قابل استقرار است. همچنین افزایش مقطع زمانی همگرایی مدل، به نسبتِ پیچیدگیِ فضای جستجو، شیبی کمتر از حالت خطی (\lr{Sublinear}) دارد.
    \item \textbf{تصمیم‌گیریِ شفاف و قابل مدل‌سازی:} رویه‌هایِ رفتارگونه‌ی عامل (همانند پدیده مصالحه در برخورد با اتلافِ دیتا و زمانِ توقف) به نحو کاملاً معنا‌دار و قابل تعریفی مدل‌گستر و قابل توضیح می‌باشد که این ویژگی به مهندسان قدرت تحلیل مضاعف و قابلیت اطمینانِ بیشتری عرضه می‌نماید.
\end{enumerate}

\section{محدودیت‌ها و تهدیدات اعتبار (\lr{Threats to Validity})}

مانند تمامی‌رهیافت‌های پژوهشی، متدولوژی معرفی‌شده با محدودیت‌هایی مواجه است که بررسی صریحِ روایی و صحتِ نتایجِ استخراج‌شده، بستر مناسب را برای پژوهش‌های مکملِ آینده تبیین می‌نماید.

\subsection{اعتبار داخلی (\lr{Internal Validity})}

\begin{enumerate}
    \item \textbf{حساسیت نسبت به ابرپارامترها:} فاکتورِ کنترل‌کننده‌ی شکل‌دهیِ تجربیات تصادفی ($\rho$)، متولی تنظیم نمودن درصدِ اختصاص تجربه‌هایِ ناموفقِ تغییرِ یافته به نمونه‌های فرضیِ مثبت و امیدوارکننده است. با اینکه از طریق کاوش شبکه‌ای جامع (\lr{Grid Search}) توانستیم مناسب‌ترین مقادیر ممکن برای بهینه‌سازیِ عملکرد یافت کنیم، لیکن این مقادیر می‌توانند به نسبت سایر سیستم‌ها و دامنه‌های متفاوت، دستخوش نوسان باشند. بدیهی‌ست چنانچه پارامتر بسیار بالا وضع گردد موجب اختلال در روند یادگیری گردیده و مقادیر ضعیفِ آن نیز سبب نقصان در عامل فشارِ اکتشاف خواهد گشت.
    \item \textbf{نوسانات بذر تصادفیِ تولید ارقام (\lr{Random Seed Sensitivity}):} با وجود اجرای بیش از ۳۰ دور فرآیندِ آموزش مستقل همراه با صحت‌سنجی مستحکم آماری، ماهیتِ درگیر با احتمالات - چه در پویایی خود محیط و چه در بدنه تنظیم پاداش‌ها - سبب می‌شود تا خروجی‌های هر دَوره آموزشی با درجاتِ کوچکی از انحراف واریانس مواجه گردند. با این وصف، فواصل و حاشیه‌های اطمینان به دست آمده (قید شده در جداول فصل ۵) تبیین‌گر کمیت این نوسانات هستند.
    \item \textbf{مفروضات بازده دودویی (\lr{Binary Reward Assumption}):} یکی از شرایط پیش‌نیازِ به کار رفته در الگوریتم، اتکا بر مفروضات مرتبط با نتایج قطعی پیامدهایِ انطباقی به‌عنوان اهدافِ توفیق/سقوطِ مطلقِ دودویی (یا در مقیاس‌های استاندارد مشابه) است که نمی‌توان آن را به سهولت با کلیه نیازهای سیستم‌های ترکیبیِ غول‌پیکر وفق داد.
\end{enumerate}

\subsection{اعتبار خارجی (\lr{External Validity})}

\begin{enumerate}
    \item \textbf{تعمیم‌پذیری ساختار و تخصیصِ حوزه (\lr{Domain Specificity}):} فرآیند تمرکز و سنجشِ نهاییِ ارزیابی این رساله منحصراً روی بستر محیط شبکه‌های متراکم و سنسورهای اینترنتی (\lr{DeltaIoT}) انجام شد. هرچند پلتفرم مذکور آینه‌ای از مشکلات اساسی سیستم‌های مدرن را به تصویر کشیده و تاییدِ کیفیت کرده است، با این وصف، اعمال الگوریتمِ توسعه‌یافته بر دیگر شبکه‌هایِ توزیع‌یافته، نیازمند اعتبارسنجیِ مستقلی خواهد بود.
    \item \textbf{تمرکز بهینه بر شبکه‌های گسسته:} تدوین و صورت‌بندی فعلی به نحوی طراحی گشته که متناسب با فضای تصمیم‌داریِ مدل‌های محدود گسسته پایه‌ریزی شود. لیکن اکثریت سیستم‌های پویا (مانند سیستم‌های مقیاس‌بندی و تنظیمِ ظرفیتِ پردازشی‌ ابری) از پارامترهای پیوسته استفاده می‌نمایند و تنظیم معماری \lr{RS-DRL} متناسب با متدهای کلاسیک پیوسته نظیر شبکه‌های \lr{DDPG} ساختارهای بسط یافته مطالعاتی جدیدی را می‌طلبد.
    \item \textbf{توسعه‌ی در حوزه‌هایِ حساسِ ایمنی:} استفاده از استراتژی خوش‌بینی بر روی اقداماتی که ذاتاً و در دنیای فیزیکی خطرساز بوده‌اند (همانند پلتفرم‌های هدایت اتومبیل‌های هوشمند و یا سیستم‌های دارورسانی)، بدون در نظر گرفتن شرط احتیاط، ممکن است باعث ترغیب ماشین به انجام فعالیت‌های منجر به حادثه گردد و افزودن حلقه‌های تکمیلی تایید صلاحیت بسیار با اهمیت است.
\end{enumerate}

\subsection{اعتبار سازه (\lr{Construct Validity})}

\begin{enumerate}
    \item \textbf{شفافیت متریک‌ها:} سنجش عملکرد عامل هوش از قبیل بررسی مجانبیِ ثبات کیفیت زمانی تا رسیدن به راندمان مطلوب آستانه، رویه‌هایی پذیرفته شده می‌باشند؛ اما معرفی فاکتورهای سنجشی جایگزین مانند نرخ رنجوری (رگرت) و شاخص پایداری ممکن است ابعاد ناشناخته دیگری را نمایان سازند.
    \item \textbf{اصطکاک سنجشی مدل‌های جایگزین:} از آنجا که مدل‌های برگرفته از \lr{HER} نیازمند تنظیم صریحِ برچسب‌گذاریِ مقصد بوده‌اند، مقایسه صورت پذیرفته نیاز به در نظر گرفتن جایگزین‌های ضمنی کیفی داشت. هرچند مستندسازی اعمال شده بسیار متوازن بود ولی باید اذعان داشت که استفاده از این مقادیرِ تعدیل یافته معادل بهینه‌ای برای تمام ظرفیت موجود این دست از رقبا نمی‌تواند قلمداد گردد.
\end{enumerate}

\section{مسیرهایی برای تحقیقات آینده}

\subsection{استراتژی‌های انطباقی و وفقی شکل‌دهی $\rho$}

همان‌طور که در فصل بررسی چالش‌ها بحث گردید، اتکا به متغیر کنترل‌گر $\rho$ برای الگوریتم ضروری است. ما در توسعه‌های آتی در نظر داریم با برنامه‌ریزی یک ساختار هوشمند برای تنظیم نرخ فاکتور کنترل‌کننده متناسب با فرآیندهای زمانی مدل پیشروی کنیم تا قابلیت استقرار در زمان اجرای سیستم را تسهیل نماییم. برخی از جهت‌گیری‌های تخصصی شامل موارد زیر می‌باشد:
\begin{enumerate}
    \item \textbf{برنامه‌ریزی تنظیمیِ پویای پارامتر ($\rho$):} به جای یک متغیر تنظیم شده‌ی دستی، از سازوکاری جهت ارتقا یا تخفیفِ متغیر بهره خواهیم جست؛ بدین گونه که در هنگامه موفقیت‌های مکرر و تسلط ماشین (فاز استثمار) از درصدِ فاکتور $\rho$ کاسته شده و در هنگام درگیری مستمرِ ماشین و گرفتار شدن آن درون تله‌ی بهینه محلی که نرخ دریافت ارزش با فُتول روبرو می‌گردد به صورت پویا با ارتقای متغیر ضریب فشارِ فرار و اکتشاف تجربیات جدید تشدید شود. در این مسیر استفاده از مدل‌های فرا-یادگیری (\lr{Meta-Learning}) نیز در خصوص استنتاج بهینه فاکتور مفید فایده خواهد بود.
    \item \textbf{شکل‌دهی هوشمند بر پایه‌ی خطای عدم‌اطمینان:} امکان به‌کارگیری ابزارهای ارزیابی تخمین اطمینان (مانند \lr{Dropout} یا انواع مدل‌های اِنسامبل) در الگوریتم، مقدور می‌سازد تا آن دسته از تجربه‌های شکست خورده‌ای در چرخه‌ی تصادفی‌سازی اولویت یابند که سیستم درباره‌ی قطعیت خطاهای‌شان کمترین اطمینان و شناخت را در محیطِ غیرشناخته از خود نشان داده بود.
    \item \textbf{اولویت‌بندی حافظه‌ی نمونه‌گیری:} ادغام روش انتخاب هوشمند با الگوریتم‌های حافظهِ مبتنی بر اولویت (\lr{Prioritized Experience Replay})، عاملی را معرفی می‌کند تا زیرمجموعه‌های مردود شده‌ی مستعدِ جبران پس از اعمال شکل‌دهی، با فرکانس شانسِ بیشتری برای بازیابی توسط شبکه مجدداً احضار شوند.
\end{enumerate}

\subsection{گسترش به دامنه‌های پیوسته و حساس به ایمنی}

با عنایت به ماهیت پویای فضای نرم‌افزار، توسعه روش‌های حاضر بر مسائل نوین‌تر ضروری خواهد بود:
\begin{enumerate}
    \item \textbf{فضای اعمال پیوسته (\lr{Continuous Action Spaces}):} بسط معماری متدیک فعلی در بَطنِ متدهایِ قدرتمندتری نظیر سیستم‌های شیبِ سیاستی همانند متدهای \lr{PPO} و \lr{SAC} از اهداف توسعه‌ی اصلی پژوهش است تا علاوه بر شکل‌دهی پاداش تصادفی، بتوان پارامترهای بدون گسستگی (همچون تنظیم منابع سخت‌افزاری پردازش در فضای ابری) را تنظیم نمود.
    \item \textbf{اعمال محدودیت‌های تایید ایمنی (\lr{Safe Reshaping}):} برای رفع نگرانی مطرح شده در دامنه‌های حساس به ایمنی، در نظر داریم مکانیزمی‌با عنوان «شکل‌دهی ایمن» پدید آوریم که بازتولیدِ تصادفیِ نتایج شکست‌خورده تنها زمانی مُیسّر گردد که خط‌مشی پیش‌بینی‌شده در مناطق ایمن تاییدشده قرار گیرد. ترکیب با ابزارهای اعتبارسنجی صوری (همچون \lr{UPPAAL-SMC} استفاده شده در این تحقیق) از هدایت به وضعیت‌های ناایمن پیشگیری می‌کند.
    \item \textbf{گسترش به سامانه‌های چند عاملی (\lr{Multi-Agent}):} برای مقابله با فضاهایی با پویاییِ بخش‌های تطبیقی در حال تعامل با هم، گسترش متد در چارچوبی صورت می‌پذیرد که چند عامل در هماهنگی استراتژی‌های کاوش همسان به شکل‌دهی پاداش‌های منسجم برای تحقق اهداف بزرگ عمل کنند.
\end{enumerate}

\subsection{اعتبارسنجی بین‌دامنه و تحلیل‌های نظری}

جهت بررسی فراگیری این استراتژی و ایجاد پایه‌های نظری قوی‌تر اهداف زیر دنبال خواهد گردید:
\begin{enumerate}
    \item \textbf{پیکر‌بندی شبکه‌های ابری (\lr{Cloud Auto-scaling}):} به‌کارگیری این رویکرد به جهت تخصیص و توزیع مقیاس ماشین‌های مجازی، و ارزیابی در برابر کارهای مرتبط دیگر در این محیط.
    \item \textbf{هماهنگی وسایل نقلیه خودران (\lr{Autonomous Vehicles}):} آزمایش در محیط‌های شبیه‌سازی برای هماهنگ‌سازی و تطبیق الگوهای ترافیکی با محدودیت‌های ایمنی و اهداف چندگانه (زمان سفر، بهینه‌سازی انرژی و امنیت).
    \item \textbf{تضمین تئوریک فرآیند همگرایی (\lr{Theoretical Analysis}):} بررسی اینکه آیا سیگنال پاداش شکل‌دهی‌شده ویژگی انقباضی و نگاشت ثابت برای یادگیری کیو (\lr{Q-Learning}) را حفظ می‌کند، و محاسبه کران‌های خطای احتمالی انحرافِ ارزش \lr{Q} نسبت به سیاست بهینه.
    \item \textbf{کاهش حجمِ پیچیدگی نمونه‌ها (\lr{Sample Complexity Analytical Analysis}):} کمیت‌بندی تئوریک پیرامون اثربخشیِ شکل‌دهی پاداش بر تقلیل شمار نمونه‌های لازم برای دستیابی به خروجی‌های نزدیک به بهینه در قیاس با شبکه‌ی \lr{DQN} استاندارد.
\end{enumerate}

\section{نتیجه‌گیری}

\subsection{خلاصه بیان مسئله و راهکار پیشنهادی}

این رساله در قامت دستاوردی متصل بر یکی از چالش‌های بنیادین سیستم‌های خودوفقی مقاصد خود را تنظیم نموده است؛ مدیریت عدم‌قطعیت در شبکه‌هایی که درون فضاهای بزرگ و گسسته‌ی تطابق عمل می‌کنند. رویکردهای سنتی - اعم از مدل‌های مبتنی بر یادگیری ماشین، مدل‌های متکی بر جستجو و یا یادگیری تقویتی استاندارد - با محدودیت‌های شدیدی مواجه بوده‌اند؛ ویژگی‌هایی از قبیل نیاز مبرم به دیتای برچسب‌دار آفلاین، دست و پنجه نرم کردن با فقدان مقیاس‌پذیری و احتمال بالای توقف در بهینه‌های محلیِ برآمده از تغییر شرایط.

ما استراتژی \textbf{یادگیری تقویتی عمیق با شکل‌دهی تصادفی پاداش (\lr{RS-DRL})} را ارائه نمودیم؛ رویکردی نوین که مکانیسم تصادفیِ شکل‌دهی را درون جریان حافظه بازپخشِ عامل کیو (\lr{DQN}) یکپارچه می‌نماید. با بهره‌گیری از زیرمجموعه‌ای از تجربیات ناکامِ فاز گذشته ($r_i=0$) و اعمال تحول در پاداش و تبدیل آن به تجربه‌ای خوش‌بینانه ($r_i=1$)، راهکار \lr{RS-DRL} مشوق قدرتمندِ فرآیندِ اکتشاف بوده و از گرفتاری در استثنائات محلی جلوگیری می‌نماید، بدون آنکه کمترین الزامی‌برای تعریف و اختصاص قطعی‌ِ اهدافِ مبتنی بر دانش از پیش‌تعیین‌شدهِ دامنه، داشته باشد.

\subsection{مرور مشارکت‌ها و شواهد تجربی}

یافته‌های مستخرجِ پژوهش که ارائه‌ی پاسخی قاطع بر پرسش‌های این رساله به شمار می‌آید شامل موارد کلیدی زیر است. جدول \ref{tab:summary_contributions} خلاصه‌ای از این مشارکت‌ها و شواهد را نشان می‌دهد.

\begin{table}[h!]
\centering
\caption{خلاصه‌ای از مشارکت‌های کلیدی و شواهد تجربی}
\label{tab:summary_contributions}
\begin{tabular}{|p{4.5cm}|p{5.5cm}|p{5cm}|}
\hline
\textbf{مشارکت (\lr{Contribution})} & \textbf{شواهد (\lr{Evidence})} & \textbf{اهمیت (\lr{Significance})} \\ \hline
شکل‌دهی پاداش نوین (\lr{Novel Reward Shaping}) & توسعه الگوریتم ۲، نمایش در شکل‌های ۵ و ۶ & امکان اکتشاف بدون نیاز به دانش دامنه سیستم \\ \hline
تطبیق و وفقی‌سازی مقیاس‌پذیر & مستخرج از جداول ۳ و ۴ (نسخه‌های \lr{v1} تا \lr{v2}) & حفظ پایداریِ عملکرد با رشد گزینه‌ها از ۲۱۶ تا ۴۰۹۶ حالت \\ \hline
برتری و تفوق اثبات آماری & تحلیل جداول ۵ و ۶، معناداری ۲۵ از ۲۸ تست & دریافت \lr{$p < 0.001$} برای اکثر متریک‌های بهبود اتلاف بسته \\ \hline
\end{tabular}
\end{table}

\textbf{مشارکت اول: ابداع ماشین نوین جهت شکل‌دهی پاداش}
توسعه الگوریتم تخصیص پاداش (طراحی الگوریتم ۲) که از طریق کنترل‌کننده‌ای با ضریب مشخصِ ($\rho$) عمل نموده و برخلاف رویکردهای غایت‌محور نظیر متدهای \lr{HER} هیچگونه نیازی بر تعیین هدف آستانه‌ای با درک معنایی نداشته و مستقیماً روی سیستم‌های خودوفقی قابل راه‌اندازی است.

\textbf{مشارکت دوم: یکپارچه‌سازی حلقه کنترل \lr{MAPE-K}}
طراحی زیرِسامانه‌ی متصل متشکل از چرخه ۹ مرحله‌ای که \lr{RS-DRL} را با حلقه \lr{MAPE-K} ترکیب می‌کرد؛ مجهز به ارزیابی برای تضمین اعتبارِ زمانِ اجرا با موتور کنترل \lr{UPPAAL-SMC} و مؤلفه \lr{ActivFORMS} برای اجرای وفقی.

\textbf{مشارکت سوم: دستاوردهای تجربی در مقیاس اینترنت اشیاء}
دستاورد استقرار شبکه‌ی عامل بر بستر ۳ مقیاس وسیع (فضای محدود ۲۱۶، نسخه ۱۲۹۶ و نسخه منبسط ۴۰۹۶ تصمیم ممکن). دستیابی بی‌رقیب به نرخ مجانبی $0.9905 \pm 0.0268$ (در شبکه اولیه) و $0.8983 \pm 0.0932$ (در \lr{v1.1}) که با قاطعیت ساختار کلاسیک \lr{$\epsilon$-greedy DQN} (با راندمان $0.4559 \pm 0.2233$) و مشابه‌های \lr{HER} را مغلوب ساخت. ثبت کاهش شگفت‌انگیزِ توالیِ ٪۲۲.۲۳ برای نابودی بسته‌ها و ٪۳۹.۰۷ کاهش تأخیر شبکه‌ نسبت به \lr{DLASER+} تحت اهداف \lr{TTS}. این مقادیر از طریق نتایج آزمون \lr{t-test} (موفقیت در ۸۹.۳٪ از ۲۸ مقایسه متمایز) اعتبار مطلق دریافت کرد.

\textbf{مشارکت چهارم: بینش عمیق نظری و عملی}
اثبات و آشکارسازیِ اینکه مکانیزمِ روایی برای تزریق خوش‌بینی نه تنها راه‌حلی منطقی، خروج سیستم را از گیر افتادن مصون داشته، بلکه تحلیل الگوهای متضاد درون‌سیستمی‌نظیر (تأخیر متناقض با نرخ اتلاف بسته‌ها) نشان دادیم عامل استنتاج‌هایی به وسعت دانش طراحی چند-هدفه و سازگار با رفتار سیستم هوشمند اخذ کرده است و همزمان این ظرفیت برای تغییر سایز و پدیده‌ی مقیاس‌پذیری فضای عمل بسیار کارآمد بوده است.

\subsection{سخنان پایانی}

از آنجا که سیستم‌های نرم‌افزاری به‌طور فزاینده‌ای در محیط‌هایی سرشار از عدم قطعیت، پویایی و مقیاس‌های بزرگ عمل می‌کنند، نیاز به مکانیسم‌های وفقی که بتوانند به‌طور کاملاً مستقل فرآیند یادگیری و تنظیم استراتژی را انجام دهند چشمگیرتر می‌شود. مدل هوشمند \lr{RS-DRL} یک قدم استوار در جهت این چشم‌انداز تلقی می‌شود - روشی که نه‌تنها متکی به آموزش‌های مبتنی بر موفقیت است بلکه با ظرافتی شایسته، از بطن اشتباهات پیشین درس‌های بسیار موثری اخذ می‌کند و به بیان دیگر کاستی‌های پیشین را به یک فرصت ارزنده تبدیل خواهد کرد.

اهمیت این تحقیقات مرزهای توسعه‌یافته در یک شبکه \lr{DeltaIoT} را کاملاً در هم شکسته است. هر نوع معماری که توأمان با ۳ چالش بنیادین سیستمِ نرم‌افزار (استقرار محیط پیکربندی کلان، عدم‌قطعیت‌های متغیر شرایط و نیاز تطبیق زمانِ اجرا) دست و پنجه نرم کند، با تکیه بر متد معرفی‌شده در \lr{RS-DRL} می‌تواند بر این معضل غلبه نماید، به گونه‌ای که سیستم‌های ابری، امنیتِ سایبریِ شبکه‌، کنترلِ برق هوشمند و هدایت پهپادهای یکپارچه از مصادیق این قابلیت‌ها هستند.

مأموریتِ نهایی آینده برای سیستم‌ها آن است که نه تنها مستمراً راهبریِ سیستم را هماهنگ بدانند بلکه فراتر از آن استراتژی را پیوسته غنی‌سازی نمایند – هدفی بزرگ که در شُرف تحقق پیوسته است. رساله‌ی حاضر با اتکا بر راهکار ابداعی \lr{RS-DRL} تلاش داشته خشتِ کوچکی اما ارزشمند رو به سوی این آرمان کلان مهندسی نرم‌افزار بنا کند.

\newpage

% کتابنامه و پیوست‌ها
\pagestyle{empty}
{
\onehalfspacing
\bibliographystyle{ieeetr-fa}
\bibliography{References}
}

\pagestyle{fancy}
\appendix
\chapter{جریان‌های کاری کامل فرآیند}
\newpage
\chapter{لیست کامل الگوریتم‌ها}
\newpage
\chapter{نتایج تجربی گسترده}
\newpage
\chapter{مصنوعات پیاده‌سازی و پکیج بازتولید}

\newpage

% چکیده و عنوان انگلیسی
\en-abstract{
(English summary of the thesis, required at the end by TMU guidelines)
}
\latinfirstPage

\end{document}
